\providecommand{\toplevelprefix}{../..}  %
\documentclass[../../book-main.tex]{subfiles}

\begin{document}

\chapter{Pursuing Low-Dimensional Distribution via Denoising}
\label{ch:general-distribution}\label{ch:compression}
\label{ch:denoising} 

\begin{quote}
\hfill  ``{\textit{Information is the resolution of uncertainty}}.''

$~$\hfill -- Claude Shannon, 1948
\end{quote}
\vspace{5mm}


In \Cref{ch:linear-independent}, we have shown how to learn simple classes of distributions whose supports are assumed to be either a single or a mixture of low-dimensional subspaces or low-rank Gaussians. For further simplicity, the different (hidden) linear or Gaussian modes are assumed to be orthogonal or independent\footnote{Or can be reduced to such idealistic cases under benign conditions.}, as illustrated in \Cref{fig:subspaces}. As we have shown, for such special distributions, one can derive rather simple and effective learning algorithms with correctness and efficiency guarantees. The geometric and statistical interpretation of operations in the associated algorithms is also very clear.

In practice, both linearity and independence are rather idealistic assumptions that distributions of real-world high-dimensional data \textit{rarely} satisfy. The only thing that we may assume is that the intrinsic dimension of the distribution is very low compared to the dimension of the ambient space in which the data are embedded. Hence, in this chapter, we show how to learn a more general class of low-dimensional distributions in a high-dimensional space that is \textit{not} necessarily (piecewise) linear.

It is typical that the distribution of real data often contains multiple components or modes, say corresponding to different classes of objects in the case of images. These modes might not be statistically independent and they may even have different intrinsic dimensions. It is also typical that we have access to only a finite number of samples of the distribution. Therefore, in  general, we may assume our data are distributed on a mixture of low-dimensional nonlinear submanifolds in a high-dimensional space. \Cref{fig:mixture-manifolds} illustrates an example of such a distribution.

To learn such a distribution under such conditions, there are several fundamental questions that we need to address:
\begin{itemize}
	\item What is a general approach to learning a general low-dimensional distribution in a high-dimensional space and representing the learned distribution?
	\item How do we measure the complexity of the distribution so that we can effectively exploit the low dimensionality for learning?
	\item How do we make the learning process computationally tractable and even scalable, given that the ambient dimension is usually high and the number of samples is typically large?
\end{itemize}
As we will see, the fundamental idea of \textit{compression}, or \textit{dimension reduction}, which has been shown to be very effective for the linear/independent case, still serves as a general principle for developing effective computational models and methods for learning general nonlinear low-dimensional distributions.

Due to its theoretical and practical significance, we will study in greater depth how this general framework of learning low-dimensional distributions via compression substantiates when the distribution of interest can be well-modeled or approximated by a mixture of low-dimensional subspaces or low-rank Gaussians. As we will see in later chapters, these seemingly restrictive classes of models serve as the basic building blocks of algorithmic development for general distributions.

\begin{figure}
    \centering
    \includegraphics[width=0.5\linewidth]{\toplevelprefix/chapters/chapter3/figs/mixed-manifolds.png}
    \caption{Data \(\vX\) are  distributed on a mixture of low-dimensional submanifolds \(\cup_{j} \cM_{j}\) in a very high-dimensional ambient space, say \(\R^{D}\). Notice that this is a much more general case than the special cases studied in the previous \Cref{ch:linear-independent} where the submanifolds were assumed to be piecewise linear hence admit parametric forms and analytical solutions.}
    \label{fig:mixture-manifolds}
\end{figure}


\section{Entropy Minimization and Compression}

\subsection{Entropy and Coding Rate}
In \Cref{ch:intro}, we have mentioned that the goal of learning is to find the simplest way to generate a given set of data. Conceptually, the Kolmogorov complexity was intended to provide such a measure of complexity, but it is not computable and is not associated with any implementable scheme that can actually reproduce the data. Hence, we need an alternative, computable, and realizable measure of complexity. That leads us to the notion of \textit{entropy}, introduced by Shannon in 1948 \cite{Shannon-1948}.

To illustrate the constructive nature of entropy, let us start with the simplest case. Suppose that we have a discrete random variable that takes \(N\) distinct values, or \textit{tokens}, \(\{\vx_{1}, \dots, \vx_{N}\}\) with equal probability \(1/N\). Then we could encode each token \(\vx_{i}\) using the \(\log_2 N\)-bit binary representation of \(i\). This coding scheme could be generalized to encoding arbitrary discrete distributions \cite{Cover-Thomas}: Given a distribution \(p\) such that \(\sum_{i=1}^{N} p(\vx_{i}) = 1\), one could assign each token \(\vx_{i}\) with probability \(p(\vx_{i})\) to a binary code of size \(\log_2 [1/p(\vx_{i})] = - \log_2 p(\vx_{i})\) bits. Hence the average number of bits, or \textit{the coding rate}, needed to encode any sample from the distribution \(p(\cdot)\) is given by the expression:\footnote{By the convention of  Information Theory \cite{Cover-Thomas}, the \(\log\) here has base \(2\). Hence, entropy is measured in (binary) bits.}
\begin{equation}
	H(\vx) \doteq \Ex_{\vx \sim p}[\log 1/p(\vx)]  = - \sum_{i=1}^{N} p(\vx_{i}) \log  p(\vx_{i}).
	\label{eqn:entropy-discrete}
\end{equation}
This is known as the \textit{entropy} of the (discrete) distribution \(p(\cdot)\). Note that this entropy is always nonnegative and it is zero if and only if \(p(\vx_{i}) = 1\) for some \(\vx_{i}\) with \(i \in [N]\).\footnote{Here notice that we use the fact \(\lim_{p\rightarrow 0} p \log p = 0\).}


\subsection{Differential Entropy}

When the random variable \(\vx \in \R^{D}\) is continuous and has a probability density \(p\), one may view that the limit of the above sum \eqref{eqn:entropy-discrete} is related to an integral:
\begin{equation}
	h(\vx) \doteq \Ex_{\vx \sim p}[\log 1/p(\vx)] = - \int_{\R^{D}} p(\vxi) \log p(\vxi) \odif{\vxi}.
	\label{eqn:entropy-differential}
\end{equation}
{More precisely, given a continuous variable \(\vx\), we may quantize it with a quantization size \(\epsilon > 0\). Denote the resulting discrete variable as \(\vx^{\epsilon}\). Then one can show that \(H(\vx^{\epsilon}) + \log(\epsilon) \approx h(\vx)\). Hence, when \(\epsilon\) is small, the differential entropy \(h(\vx)\) can be negative. Interested readers may refer to \cite{Cover-Thomas} for a more detailed explanation.}

\begin{example}[Entropy of Gaussian Distributions]
	Through direct calculation, it is possible to show that the entropy of a Gaussian distribution \(\x \sim \dNorm(\mu, \sigma^2)\) is given by:
	\begin{equation}
		h(x) = \frac{1}{2}\log (2\pi \sigma^{2}) + \frac{1}{2}.
		\label{eqn:entropy-Gaussian}
	\end{equation}
	It is also known (cf. \Cref{thm:max_entropy}) that the Gaussian distribution achieves the maximal entropy for all distributions with the same variance \(\sigma^{2}\). The entropy of a multivariate Gaussian distribution \(\vx \sim \dNorm(\vmu, \vSigma)\) in \(\R^{D}\) is given by:
	\begin{equation}
		h(\vx) = \frac{D}{2}(1 + \log(2\pi)) + \frac{1}{2}\logdet(\vSigma).
		\label{eqn:entropy-Gaussian-multi}
	\end{equation}
\end{example}

Similar to the entropy of a discrete distribution, we would like the differential entropy to be associated with the coding rate of some realizable coding scheme. For example, as mentioned above, we may discretize the domain of the distribution with a grid of size \(\epsilon > 0\). The coding rate of the resulting discrete distribution can be viewed as an approximation to the differential entropy \cite{Cover-Thomas}.

Be aware that there are some caveats associated with the definition of differential entropy. For a distribution in a high-dimensional space, when its support becomes degenerate (low-dimensional), its differential entropy diverges to \(-\infty\). This fact is proved in \Cref{thm:max_entropy} (we also recall the maximum entropy characterization of the Gaussian distribution mentioned above in \Cref{thm:max_entropy}) but even in the simple explicit case of Gaussian distributions \eqref{eqn:entropy-Gaussian-multi}, when the covariance \(\vSigma\) is singular, we can see that \(\logdet(\vSigma) = -\infty\) so we have \(h(\vx) = -\infty\). In such a situation, it is not obvious how to properly quantize or encode such a distribution. Nevertheless, degenerate (Gaussian) distributions are precisely the simplest possible, and arguably the most important instances of low-dimensional distributions in a high-dimensional space. In this chapter, we will discuss a complete resolution to this seeming difficulty with degeneracy.


\subsection{Minimizing Coding Rate}\label{sub:min_entropy}
Remember that the learning problem entails the recovery of a (potentially continuous)  distribution \(p(\vx)\) from a set of samples \(\{\vx_{1}, \ldots, \vx_{N}\}\) drawn from the distribution. For ease of exposition, we write \(\vX = [\vx_{1}, \ldots, \vx_{N}] \in \R^{D\times N}\) . Given that the distributions of interest here are (nearly) low-dimensional, we should expect that their (differential) entropy is very small. However, unlike the situations we studied in the previous chapter, in general, we d,o not know the family of (analytical) low-dimensional models to which the distribution \(p(\vx)\) belongs. So checking whether the entropy is small seems to be the only guideline that we can rely on to identify and model the distribution.

Now given the samples alone without knowing what \(p(\vx)\) is, in theory they could be interpreted as samples from any generic distribution. In particular, they could be interpreted as any of the following:
\begin{enumerate}
	\item as samples from the empirical distribution \(p^{\vX}\) itself, which assigns \(1/N\) probability to each of the \(N\) samples \(\vx_{i}, i=1, \ldots, N\).
	\item as samples from a standard normal distribution \(\vx^{n} \sim p^{n} \doteq \dNorm(\vzero, \sigma^{2} \vI)\) with a variance \(\sigma^{2}\) large enough (say larger than the sample norms);
	\item as samples from a normal distribution \(\vx^{e} \sim p^{e} \doteq \dNorm(\vzero, \hat{\vSigma})\) with a covariance \(\hat{\vSigma} = \frac{1}{N} \vX \vX^{\top}\) being the empirical covariance of the samples;
	\item as samples from a distribution \(\hat \vx \sim \hat{q}(\vx)\) that closely approximates the ground truth  distribution \(p\).
\end{enumerate}
Now the question is: which case is better, and in what sense? Suppose that you believe these data \(\vX\) are drawn from a particular distribution  \(q(\vx)\), which may be one of the above distributions considered. Then we could encode the data points with the optimal code book for the distribution \(q(\vx)\). The required average coding length (or coding rate) is given by the so-called \textit{cross-entropy}:
\begin{equation}
	\CE(p \mmid q) \doteq \Ex_{\vx \sim p}[\log 1/q(\vx)],
\end{equation}
as the number of samples \(N\) becomes large. If we have identified the correct distribution \(p(\vx)\), the coding rate is given by the entropy \(h(\vx) = - \Ex_{\vx \sim p}[\log 1/p(\vx)]\). It turns out that the above coding length \(\CE(p \mmid q)\) is always larger than or equal to the entropy \(h(\vx)\) unless \(p = q\). 

Hence, given a set of sampled data \(\vX\), to determine which distribution is a better characterization of the data among \(p^{n}\), \(p^{e}\),  and \(\hat{q}\), we may compare their coding rates for \(\vX\) and see which one gives the lowest rate. We know from the above that the (theoretically achievable) coding rate for a distribution is closely related to its entropy. In general, we have:
\begin{equation}
	h(p^{n}) > h(p^{e}) > h(\hat{q}).
\end{equation}
Hence, if the data \(\vX\) were encoded by the code book associated with each of these distributions, the coding rate for \(\vX\) would in general decrease in the same order:
\begin{equation}
	\CE(p \mmid p^{n}) > \CE(p \mmid p^{e}) > \CE(p \mmid \hat{q}).
\end{equation}

This observation gives us a general guideline on how we may be able to pursue a distribution \(p(\vx)\) which has a low-dimensional structure. It suggests two possible approaches:
\begin{enumerate}
	\item Starting with a general distribution (say a normal distribution) with high entropy, gradually transform the distribution towards the (empirical) distribution of the data by reducing entropy.
	\item Among a large family of (parametric or non-parametric) distributions with explicit coding schemes that encode the given data, progressively search for better coding schemes that give lower coding rates.
\end{enumerate}
Conceptually, both approaches essentially try to do the same thing. For the first approach, we need to make sure such a path of transformation exists and is computable. For the second approach, it is necessary that the chosen family is rich enough and can closely approximate (or contain) the ground truth distribution. For either approach, we need to ensure that solutions with lower entropy or better coding rates can be efficiently computed and converge to the desired distribution quickly.\footnote{Say the distribution of real-world data such as images and texts.} We will explore both approaches in the two remaining sections of this chapter.  %

\section{Compression via Denoising}\label{sub:compression_denoising}

In this section, we will describe a \textit{natural} and \textit{computationally tractable} way to learn a distribution \(p(\vx)\) by way of learning a parametric encoding of our distribution such that the representation has the minimum entropy or coding rate, then using this encoding to transform high-entropy samples from a standard Gaussian into low-entropy samples from the target distribution, as illustrated in \Cref{fig:diffusion-chapter3}. This presents a methodology that utilizes both approaches above in order to learn and sample from the distribution.

\begin{figure}[t]
	\centering
	\includegraphics[width=\linewidth]{\toplevelprefix/chapters/chapter3/figs/diffusion_pipeline.png}
	\caption{Illustration of an iterative denoising process that, starting from an isotropic Gaussian distribution, converges to an arbitrary data distribution. }
	\label{fig:diffusion-chapter3}
\end{figure}


\subsection{Diffusion and Denoising Processes} \label{sub:intro_diffusion_denoising}

We first want to find a procedure to decrease the entropy of a given very noisy sample into a lower-entropy sample from the data distribution. Here, we describe a potential approach—one of many, but perhaps the most natural way to attack this problem. First, we find a way to \textit{gradually increase} the entropy of existing samples from the data distribution. Then, we find an \textit{approximate inverse} of this process. But in general, the operation of increasing entropy does not have an inverse, as information from the original distribution may be destroyed. We will thus tackle a special case where (1) the operation of adding entropy takes on a simple, computable, and reversible form; (2) we can obtain a (parametric) encoding of the data distribution, as alluded to in the above pair of approaches. As we will see, the above two factors will ensure that our approach is possible.

\subsubsection{Diffusion Processes}

We will increase the entropy in arguably the simplest possible way, i.e., \textit{adding isotropic Gaussian noise}. More precisely, given the random variable \(\vx\), we can consider the \textit{stochastic process} \((\vx_{t})_{t \in [0, T]}\) which adds gradual noise to it, i.e.,
\begin{equation}\label{eq:additive_gaussian_noise_model}
	\vx_{t} \doteq \vx + t\vg, \qquad \forall t \in [0, T],
\end{equation}
where \(T \in [0, \infty)\) is a time horizon and \(\vg \sim \dNorm(\vzero, \vI)\) is drawn independently of \(\vx\). This process is an example of a \textit{diffusion process}, so-named because it spreads the probability mass out over all of \(\R^{D}\) as time goes on, increasing the entropy over time. This intuition is confirmed graphically by \Cref{fig:ve_forward_density}, and rigorously via the following theorem.
\begin{theorem}[Simplified Version of \Cref{thm:diffusion_entropy_increases}]
	Suppose that \((\vx_{t})_{t \in [0, T]}\) follows the model \eqref{eq:additive_gaussian_noise_model}. For any \(t \in (0, T]\), the random variable \(\vx_{t}\) has differential entropy \(h(\vx_{t}) > -\infty\). Moreover, under certain technical conditions on \(\vx\), 
	\begin{equation}
		\odv*{h(\vx_{t})}{t} > 0, \qquad \forall t \in (0, T],
	\end{equation}
	showing that the entropy of the noised \(\vx\) increases over time \(t\).
\end{theorem}
The proof is elementary, but it is rather long, so we postpone it to \Cref{sub:diffusion_entropy_increases}. The main as-yet unstated implication of this result is that \(h(\vx_{t}) > h(\vx)\) for every \(t > 0\). To see this, note that if \(h(\vx) = -\infty\) then \(h(\vx_{t}) > -\infty\) for all \(t > 0\),  and if \(h(\vx) > -\infty\) then \(h(\vx_{t}) = h(\vx) + \int_{0}^{t}[\odv*{h(\vx_{s})}{s}]\odif{s} > h(\vx)\) by the fundamental theorem of calculus, so in both cases \(h(\vx_{t}) > h(\vx)\) for every \(t > 0\).

\begin{figure}
	\includegraphics[width=\textwidth]{\toplevelprefix/chapters/chapter3/figs/ve_forward_diffusion_density.png}
	\caption{\small\textbf{Diffusing a mixture of Gaussians \(\vx\).} From left to right, we observe the evolution of the density \(p_{t}\) of \(\vx_{t}\) as \(t\) grows from \(0\) to \(10\), along with some representative samples (red). The plane is colored by the probability density \(p_{t}\); high-density regions are colored darker blue. We observe that the probability mass becomes less concentrated as \(t\) increases, signaling that entropy increases.}
	\label{fig:ve_forward_density}
\end{figure}

\begin{example}[Adding Noise to a Mixture of Gaussians]\label{ex:gaussian_mixture}
    In this example, we consider the process of adding noise to a mixture of \(K\) Gaussians. Namely, let us choose \(K\) means \(\vmu_{1}, \dots, \vmu_{K} \in \R^{D}\), \(K\) covariance matrices \(\vSigma_{1}, \dots, \vSigma_{K} \in \PSD(D)\), and probabilities \(\pi_{1}, \dots, \pi_{K} \in [0, 1]\) such that \(\sum_{k = 1}^{K}\pi_{k} = 1\). Then our data \(\vx\) is sampled as follows. 
	\begin{itemize}
		\item First, an index (or \textit{label}) \(y \in [K]\) is sampled such that \(y = k\) with probability \(\pi_{k}\).
		\item Second, \(\vx\) is sampled from the normal distribution \(\dNorm(\vmu_{y}, \vSigma_{y})\). Namely, if \(y = k\) then we sample the value of \(\vx\) from \(\dNorm(\vmu_{k}, \vSigma_{k})\).
	\end{itemize}
	Then \(\vx\) has a Gaussian mixture model distribution
	\begin{equation}
		\vx \sim \sum_{k = 1}^{K}\pi_{k}\dNorm(\vmu_{k}, \vSigma_{k}).
	\end{equation}
	Gaussian mixture models are relevant because they can approximate many types of distributions: in particular, given a distribution for \(\vx\), there is a Gaussian mixture model that can approximate it arbitrarily well. For this reason, Gaussian mixture models are a popular analytically tractable surrogate for multi-modal and complex data distributions. In this book, we will use them extensively as a ``model problem'' to understand and develop almost every concept. For now, we use a Gaussian mixture model to represent the data \(\vx\). Upon adding noise, we can see that the distribution of \(\vx_{t}\) is also a Gaussian mixture model:
	\begin{equation}
		\vx_{t} = \vx + t\vg \sim \sum_{k = 1}^{K}\pi_{k}\dNorm(\vmu_{k}, \vSigma_{k} + t^{2}\vI).
	\end{equation}
	To understand through example how the entropy changes over time, suppose that \(K=1\), i.e., \(\vx \sim \dNorm(\vmu, \vSigma)\). It follows from \eqref{eqn:entropy-Gaussian-multi} that the entropy of \(\vx_{t}\) is 
	\begin{equation}
		h(\vx_{t}) = \frac{D}{2}(1 + \log(2\pi)) + \frac{1}{2}\logdet(\vSigma + t^{2}\vI). 
	\end{equation}
	Then, we can calculate the time derivative of the entropy directly (see \Cref{app:optimization}), obtaining
	\begin{equation}
		\odv*{h(\vx_{t})}{t} = t\tr((\vSigma + t^{2}\vI)^{-1}) = \sum_{i = 1}^{D}\frac{t}{\lambda_{i} + t^{2}} > 0
	\end{equation}
	where \(\lambda_{i}\) are the (non-negative) eigenvalues of \(\vSigma\). Since this derivative is positive, the entropy \(h(\vx_{t})\) increases over time.
\end{example}


\subsubsection{Denoising Processes}
The inverse operation to adding noise is known as \textit{denoising}. It is a classical and well-studied topic in signal processing and system theory, such as the Wiener filter and the Kalman filter (see \Cref{sec:analytical}). Several problems discussed in \Cref{ch:classic}, such as PCA, ICA, and Dictionary Learning, are specific instances of the denoising problem. For a fixed \(t\) and the additive Gaussian noise model \eqref{eq:additive_gaussian_noise_model}, the denoising problem can be formulated as attempting to learn a function \(\bar{\vx}^{\ast}(t, \cdot)\) which forms the best possible approximation (in expectation) of the true random variable \(\vx\), given both \(t\) and \(\vx_{t}\):
\begin{equation}\label{eq:denoising_loss}
	\bar{\vx}^{\ast}(t, \cdot) \in \argmin_{\bar{\vx}(t, \cdot)}\Ex_{\vx, \vx_{t}}\norm{\vx - \bar{\vx}(t, \vx_{t})}_{2}^{2}.
\end{equation}
The solution to this problem, when optimizing \(\bar{\vx}(t, \cdot)\) over all possible (square-integrable) functions, is the so-called \textit{Bayes optimal denoiser}: 
\begin{equation}\label{eq:optimal_denoiser}
	\bar{\vx}^{\ast}(t, \vxi) \doteq \Ex[\vx \mid \vx_{t} = \vxi].
\end{equation}

This expression justifies the notation \(\bar{\vx}\), which is meant to compute a conditional expectation (i.e., conditional mean or conditional average). In short, it attempts to remove the noise from the noisy input, outputting the best possible guess (in expectation and w.r.t.~the \(\ell^{2}\)-distance) of the (de-noised) original random variable.

\begin{example}[Denoising Gaussian Noise from a Mixture of Gaussians]\label{example:denoising_gaussian_mixture}
	In this example we compute the Bayes optimal denoiser for an incredibly important class of distributions, the Gaussian mixture model in \Cref{ex:gaussian_mixture}. Let us define \(\phi(\vx ; \vmu, \vSigma)\) as the probability density of \(\dNorm(\vmu, \vSigma)\) evaluated at \(\vx\). In this notation, the density of \(\vx_{t}\) is 
	\begin{equation}
		p_{t}(\vx_{t}) = \sum_{k = 1}^{K}\pi_{k}\phi(\vx_{t} ; \vmu_{k}, \vSigma_{k} + t^{2}\vI).
	\end{equation}

	Conditioned on \(y\), the variables are jointly Gaussian: if we say that \(\vx = \vmu_{y} + \vSigma_{y}^{1/2}\vu\) where \((\cdot)^{1/2}\) is the matrix square root and \(\vu \sim \dNorm(\vzero, \vI)\) independently of \(y\) (and \(\vg\)), then we have
	\begin{equation}
		\mat{\vx \\ \vx_{t}} = \mat{\vmu_{y} \\ \vmu_{y}} + \mat{\vSigma_{y}^{1/2} & \vzero \\ \vSigma_{y}^{1/2} & t\vI}\mat{\vu \\ \vg}.
	\end{equation}
	This shows that \(\vx\) and \(\vx_{t}\) are jointly Gaussian (conditioned on \(y\)) as claimed. Thus we can write
	\begin{equation}
		\mat{\vx \\ \vx_{t}} \sim \dNorm\rp{\mat{\vmu_{y} \\ \vmu_{y}}, \mat{\vSigma_{y} & \vSigma_{y} \\ \vSigma_{y} & \vSigma_{y} + t^{2}\vI}}.
	\end{equation}
	Thus the conditional expectation of \(\vx\) given \(\vx_{t}\) (i.e., the
	Bayes optimal denoiser conditioned on \(y\)) is famously
	(\Cref{exercise:conditional_gaussian})
	\begin{equation}
		\Ex[\vx \mid \vx_{t}, y] = \vmu_{y} + \vSigma_{y}(\vSigma_{y} + t^{2}\vI)^{-1}(\vx_{t} - \vmu_{y}).
	\end{equation}
	To find the overall Bayes optimal denoiser, we use the law of iterated expectation, obtaining
	\begin{align}
		\bar{\vx}^{\ast}(t, \vx_{t})
		&= \Ex[\vx \mid \vx_{t}] \\ 
		&= \Ex[\Ex[\vx \mid \vx_{t}, y] \mid \vx_{t}] \\ 
		&= \sum_{k = 1}^{K}\Pr[y = k \mid \vx_{t}]\Ex[\vx \mid \vx_{t}, y = k].
	\end{align}
	The probability can be dealt with as follows. Let \(p_{t \mid y}\) be the probability density of \(\vx_{t}\) conditioned on the value of \(y\). Then
	\begin{align}
		\Pr[y = k \mid \vx_{t}]
		&= \frac{p_{t \mid y}(\vx_{t} \mid k)\pi_{k}}{p_{t}(\vx_{t})} \\ 
		&= \frac{\pi_{k}\phi(\vx_{t} ; \vmu_{k}, \vSigma_{k}
		+ t^{2}\vI)}{\sum_{i = 1}^{K}\pi_{i}\phi(\vx_{t} ; \vmu_{i}, \vSigma_{i} + t^{2}\vI)}.
	\end{align}
	On the other hand, the conditional expectation is as described before:
	\begin{equation}
		\Ex[\vx \mid \vx_{t}, y = k] = \vmu_{k} + \vSigma_{k}(\vSigma_{k} + t^{2}\vI)^{-1}(\vx_{t} - \vmu_{k}).
	\end{equation}
	So putting this all together, the true Bayes optimal denoiser is 
	\begin{equation}\label{eq:gmm_bayes_optimal_denoiser}
		\bar{\vx}^{\ast}(t, \vx_{t}) = \sum_{k = 1}^{K}\frac{\pi_{k}\phi(\vx_{t}
		;\vmu_{k}, \vSigma_{k} + t^{2}\vI)}{\sum_{i = 1}^{K}\pi_{i}\phi(\vx_{t}
		;\vmu_{i}, \vSigma_{i} + t^{2}\vI)}\cdot\bp{\vmu_{k} + \vSigma_{k}(\vSigma_{k} + t^{2}\vI)^{-1}(\vx_{t} - \vmu_{k})}.
	\end{equation}
	This example is particularly important, and several special cases will give us great conceptual insight later. For now, let us attempt to extract some geometric intuition from the functional form of the optimal denoiser \eqref{eq:gmm_bayes_optimal_denoiser}.

	To try to understand \eqref{eq:gmm_bayes_optimal_denoiser} intuitively, let us first set \(K = 1\) (i.e., one Gaussian) such that \(\vx \sim \dNorm(\vmu, \vSigma)\). Let us then diagonalize \(\vSigma = \vV\vLambda \vV^{\top}\). Then the Bayes optimal denoiser is 
	\begin{align}
		\bar{\vx}^{\ast}(t, \vx_{t}) & = \vmu + \vSigma(\vSigma + t^{2}\vI)^{-1}(\vx_{t} - \vmu) \\
        & = \vmu + \vV\mat{\lambda_{1}/(\lambda_{1} + t^{2}) & & \\ & \ddots & \\ & & \lambda_{D}/(\lambda_{D} + t^{2})}\vV^{\top}(\vx_{t} - \vmu),\label{eq:gaussian_denoiser}
	\end{align}
	where \(\lambda_{1}, \dots, \lambda_{D}\) are the eigenvalues of \(\vSigma\). We can observe that this denoiser has three steps:
	\begin{itemize}
		\item Translate the input \(\vx_{t}\) by \(\vmu\).
		\item Contract the (translated) input \(\vx_{t} - \vmu\) in each eigenvector direction by a quantity \(\lambda_{i}/(\lambda_{i} + t^{2})\). If the translated input is low-rank and some eigenvalues of \(\vSigma\) are zero, these directions get immediately contracted to \(0\) by the denoiser, ensuring that the output of the contraction is similarly low-rank.
		\item Translate the output back by \(\vmu\).
	\end{itemize}
	This is the geometric interpretation of the denoiser of a \textit{single} Gaussian. The overall denoiser of the Gaussian mixture model \eqref{eq:gmm_bayes_optimal_denoiser} uses \(K\) such denoisers, weighting their output by the posterior probabilities \(\Pr[y = k \mid \vx_{t}]\). If the means of the Gaussians are well-separated, these posterior probabilities are very close to \(0\) or \(1\) near each mean or cluster. In this regime, the overall denoiser \eqref{eq:gmm_bayes_optimal_denoiser} has the same geometric interpretation as the above single Gaussian denoiser.
\end{example}


\begin{figure}
	\centering 
	\includegraphics[width=\textwidth]{\toplevelprefix/chapters/chapter3/figs/ve_forward_diffusion_denoising.png}\vspace{-0.15in}
	\caption{\small \textbf{Bayes optimal denoiser and score of a Gaussian mixture model.} In the same setting as \Cref{fig:ve_forward_density}, we demonstrate the effect of the Bayes optimal denoiser \(\bar{\vx}^{\ast}\) by plotting \(\vx_{t}\) (red) and \(\bar{\vx}^{\ast}(t, \vx_{t})\) (green) for some choice \(t\) and \(\vx_{t}\). By Tweedie's formula \Cref{thm:tweedie}, the residual between them is proportional to the so-called (Hyv\"arinen) score \(\nabla_{\vx_{t}}\log p_{t}(\vx_{t})\). We can see that the score points towards the modes of the distribution of \(\vx_{t}\).}
	\label{fig:ve_forward_denoising}
\end{figure}

Intuitively, and as we can see from \Cref{example:denoising_gaussian_mixture}, the Bayes optimal denoiser \(\bar{\vx}^{\ast}(t, \cdot)\) should move its input \(\vx_{t}\) towards the modes of the distribution of \(\vx\). It turns out that, actually, we can quantify this by showing that the Bayes optimal denoiser \textit{takes a gradient ascent step} on the (log-)density of \(\vx_{t}\), which (recall) we denoted \(p_{t}\). That is, following the denoiser means moving from the input iterate to a region of higher probability within this (perturbed) distribution. For small \(t\), the perturbation is small so our initial intutition is therefore (almost) exactly right. The picture is visualized in \Cref{fig:ve_forward_denoising} and rigorously formulated as Tweedie's formula \cite{Robbins1956AnEB}.
\begin{theorem}[Tweedie's Formula]\label{thm:tweedie}
	Suppose that \((\vx_{t})_{t \in [0, T]}\) obeys \eqref{eq:additive_gaussian_noise_model}. Let \(p_{t}\) be the density of \(\vx_{t}\) (as previously declared). Then 
	\begin{equation}\label{eq:tweedie}
		\Ex[\vx \mid \vx_{t}] = \vx_{t} + t^{2}\nabla_{\vx_{t}} \log p_{t}(\vx_{t}).
	\end{equation}
\end{theorem}
\begin{proof}
	For the proof let us suppose that \(\vx\) has a density (even though the
	theorem is true without this assumption), and call this density \(p\). Let
	\(p_{0 \mid t}\) and \(p_{t \mid 0}\) be the conditional densities of \(\vx
	= \vx_{0}\) given \(\vx_{t}\) and \(\vx_{t}\) given \(\vx\) respectively.
	Let \(\phi(\vx ; \vmu, \vSigma)\) be the density of \(\dNorm(\vmu,
	\vSigma)\) evaluated at \(\vx\), so that \(p_{t \mid 0}(\vx_{t} \mid \vx)
	= \phi(\vx_{t} ; \vx, t^{2}\vI)\). Then a simple calculation gives 
	\begin{align}
		\nabla_{\vx_{t}} \log p_{t}(\vx_{t})
		&= \frac{\nabla_{\vx_{t}}p_{t}(\vx_{t})}{p_{t}(\vx_{t})} \\
		&= \frac{1}{p_{t}(\vx_{t})}\nabla_{\vx_{t}}\int_{\R^{D}}p(\vx)p_{t \mid 0}(\vx_{t} \mid \vx)\odif{\vx} \\
		&=
		\frac{1}{p_{t}(\vx_{t})}\nabla_{\vx_{t}}\int_{\R^{D}}p(\vx)\phi(\vx_{t}
		;\vx, t^{2}\vI)\odif{\vx} \\
		&=
		\frac{1}{p_{t}(\vx_{t})}\int_{\R^{D}}p(\vx)[\nabla_{\vx_{t}}\phi(\vx_{t}
		;\vx, t^{2}\vI)]\odif{\vx} \\
		&= \frac{1}{p_{t}(\vx_{t})}\int_{\R^{D}}p(\vx)\phi(\vx_{t} ; \vx, t^{2}\vI)\bs{-\frac{\vx_{t} - \vx}{t^{2}}}\odif{\vx} \\
		&= \frac{1}{t^{2}p_{t}(\vx_{t})}\int_{\R^{D}}p(\vx)\phi(\vx_{t} ; \vx, t^{2}\vI)[\vx - \vx_{t}]\odif{\vx} \\
		&= \frac{1}{t^{2}p_{t}(\vx_{t})}\int_{\R^{D}}p(\vx)\phi(\vx_{t} ; \vx,
		t^{2}\vI)\vx \odif{\vx}\nonumber \\
		& \quad - \frac{\vx_{t}}{t^{2}p_{t}(\vx_{t})}\int_{\R^{D}}p(\vx)\phi(\vx_{t} ; \vx, t^{2}\vI)\odif{\vx} \\
		&= \frac{1}{t^{2}p_{t}(\vx_{t})}\int_{\R^{D}}p(\vx)p_{t \mid 0}(\vx_{t} \mid \vx)\vx \odif{\vx} - \frac{\vx_{t}}{t^{2}p_{t}(\vx_{t})}p_{t}(\vx_{t}) \\
		&= \frac{1}{t^{2}p_{t}(\vx_{t})}\int_{\R^{D}}p_{t}(\vx_{t})p_{0 \mid t}(\vx \mid \vx_{t})\vx \odif{\vx} - \frac{\vx_{t}}{t^{2}p_{t}(\vx_{t})}p_{t}(\vx_{t}) \\
		&= \frac{1}{t^{2}}\int_{\R^{D}}p_{0 \mid t}(\vx \mid \vx_{t})\vx \odif{\vx} - \frac{\vx_{t}}{t^{2}} \\
		&= \frac{1}{t^{2}}\Ex[\vx \mid \vx_{t}] - \frac{\vx_{t}}{t^{2}} \\
		&= \frac{\Ex[\vx \mid \vx_{t}] - \vx_{t}}{t^{2}}.
	\end{align}
	Simple rearranging of the above equality proves the theorem.
\end{proof}
This result develops a connection between denoising and optimization: the Bayes-optimal denoiser takes a single step of gradient ascent on the perturbed data density \(p_{t}\), and the step size adaptively becomes smaller (i.e., taking more precise steps) as the perturbation to the data distribution grows smaller. The quantity \(\nabla_{\vx_{t}}\log p_{t}(\vx_{t})\) is called the \textit{(Hyv\"arinen) score} and frequently appears in discussions about denoising, etc.; it first appeared in a paper of Aapo Hyv\"arinen in the context of ICA \cite{hyvarinen05a}.

Similar to how one step of gradient descent is almost never sufficient to
minimize an objective in practice when initializing far from the optimum, the
output of the Bayes-optimal denoiser \(\bar{\vx}^{\ast}(t, \cdot)\) is almost never contained in a high-probability region of the data distribution when \(t\) is large, \textit{especially} when the data have low-dimensional structures. We illustrate this point explicitly in the following example.
\begin{example}[Denoising a Two-Point Mixture]\label{example:denoising_twopoints}
	Let \(x\) be uniform on the two-point set \(\{-1, +1\}\) and let \((x_{t})_{t \in [0, T]}\) follow \eqref{eq:additive_gaussian_noise_model}. This is precisely a degenerate Gaussian mixture model with priors equal to \(\frac{1}{2}\), means \(\{-1, +1\}\), and covariances both equal to \(0\). For a fixed \(t > 0\) we can use the calculation of the Bayes-optimal denoiser in \eqref{eq:gmm_bayes_optimal_denoiser} to obtain (proof as exercise)
	\begin{equation}
		\bar{x}^{\ast}(t, x_{t}) = \frac{\phi(x_{t} ; +1, t^{2}) - \phi(x_{t}
		; -1, t^{2})}{\phi(x_{t} ; 1, t^{2}) + \phi(x_{t} ; -1, t^{2})} = \tanh\rp{-\frac{x_{t}}{t^{2}}}.
	\end{equation}
	For \(t\) near \(0\), this quantity is near \(\{-1, +1\}\) for almost all inputs \(\bar{x}^{\ast}(t, x_{t})\). However, for \(t\) large, this quantity is not necessarily even approximately in the original support of \(x\), which, remember, is \(\{-1, +1\}\). In particular, for \(x_{t} \approx 0\) it holds \(\bar{x}^{\ast}(t, x_{t}) \approx 0\) which lies completely in between the two possible points. Thus \(\bar{x}^{\ast}\) \textit{will not output ``realistic'' \(x\)}. Or more mathematically, the distribution of \(\bar{x}(t, x_{t})\) is very different from the distribution of \(x\).
\end{example}


\subsubsection{Incremental Denoising.}
Therefore, if we want to denoise the very noisy sample \(\vx_{T}\) (where—recall—\(T\) is the maximum time), we cannot just use the denoiser \textit{once}. Instead, we must use the denoiser many times, analogously to gradient descent with \textit{decaying step sizes}, to converge to a stationary point \(\hat{\vx}\). Namely, we shall use the denoiser to go from \(\vx_{T}\) to \(\hat{\vx}_{T - \delta}\) which approximates \(\vx_{T - \delta}\), then from \(\hat{\vx}_{T - \delta}\) to \(\hat{\vx}_{T - 2\delta}\), etc., all the way from \(\hat{\vx}_{\delta}\) to \(\hat{\vx} = \hat{\vx}_{0}\). Each time we take a denoising step, the action of the denoiser becomes more like a gradient step on the original (log-)density. 

More formally, we uniformly discretize \([0, T]\) into \(L + 1\) timesteps \(0 = t_{0} < t_{1} < \cdots < t_{L} = T\), i.e.,
\begin{equation}
	t_{\ell} = \frac{\ell}{L}T, \qquad \ell \in \{0, 1, \dots, L\}.
\end{equation}
Then for each \(\ell \in [L] = \{1, 2, \dots, L\}\), going from \(\ell = L\) to \(\ell = 1\), we can run the iteration
\begin{align}
	\hat{\vx}_{t_{\ell - 1}}
	&= \Ex[\vx_{t_{\ell - 1}} \mid \vx_{t_{\ell}} = \hat{\vx}_{t_{\ell}}] \\
	&= \Ex[\vx + t_{\ell - 1}\vg \mid \vx_{t_{\ell}} = \hat{\vx}_{t_{\ell}}] \\
	&= \Ex\rs{\vx + t_{\ell - 1}\cdot\frac{\vx_{t_{\ell}} - \vx}{t_{\ell}} \given \vx_{t_{\ell}} = \hat{\vx}_{t_{\ell}}} \\
	&= \frac{t_{\ell - 1}}{t_{\ell}}\hat{\vx}_{t_{\ell}} + \bp{1 - \frac{t_{\ell - 1}}{t_{\ell}}}\Ex[\vx \mid \vx_{t_{\ell}} = \hat{\vx}_{t_{\ell}}] \\
	&= \bp{1 - \frac{1}{\ell}} \hat{\vx}_{t_{\ell}} + \frac{1}{\ell} \bar{\vx}^{\ast}(t_{\ell}, \hat{\vx}_{t_{\ell}}). \label{eq:iter denoise}
\end{align}

At the beginning of the iteration, where \(\ell\) is large, we barely trust the output of the denoiser and mostly keep the current iterate. This makes sense, as the denoiser can have huge variance (cf. \Cref{example:denoising_twopoints}). When \(\ell\) is small, the denoiser will ``lock on'' to the modes of the data distribution, as a denoising step basically takes a gradient step on the true distribution's log-density, and we can trust it not to produce unreasonable samples, so the denoising step mostly involves the output of the denoiser. At \(\ell = 1\) we even throw away the current iterate and just keep the output of the denoiser.

\begin{figure}[t]
	\centering
	\includegraphics[width=\textwidth]{\toplevelprefix/chapters/chapter3/figs/ve_gmm_denoising.png}
	\caption{\small\textbf{Denoising a low-rank mixture of Gaussians.} Each figure represents samples from the true data distribution (gray, orange, red) and samples undergoing the denoising process \eqref{eq:denoising-iteration-basic} (light blue). At top left, the process has just started, and the noise is very large. As the process continues, the noise is pushed further towards the support of the low-rank data distribution. Finally, in the bottom right, the generated samples are perfectly aligned with the support of the data and look very much like samples drawn from the low-rank Gaussian mixture model.}
	\label{fig:ve_gmm_denoising}
\end{figure}

We visualize the convergence process in \(\R^{3}\) in \Cref{fig:ve_gmm_denoising}. While we will develop some rigorous results about convergence later, we now give a closed-form example of the iterative denoising process applied to denoising a Gaussian data distribution.

\begin{example}[Iterative Denoising for a Single Gaussian] \label{ex:single_gaussian_denoising}
    To understand the above iterative denoising iteration intuitively, we apply it to denoise \(\vx_{T}\) defined in \eqref{eq:additive_gaussian_noise_model}, where \(\vx \sim \dNorm(\vmu, \vSigma)\). Using \eqref{eq:iter denoise}, we compute
	\begin{align}
		\hat{\vx}_{t_{\ell - 1}} - \vmu 
		&= \bp{1 - \frac{1}{\ell}}\hat{\vx}_{t_{\ell}} + \frac{1}{\ell}\bar{\vx}^{\ast}(t_{\ell}, \hat{\vx}_{t_{\ell}}) - \vmu \\
		&= \bp{1 - \frac{1}{\ell}}(\hat{\vx}_{t_{\ell}} - \vmu) + \frac{1}{\ell}(\bar{\vx}^{\ast}(t_{\ell}, \hat{\vx}_{t_{\ell}}) - \vmu).
	\end{align}
	Then, plugging in the form for \(\bar{\vx}^{\ast}\) that we obtain from \eqref{eq:gaussian_denoiser} and performing some tedious matrix algebra, we obtain
	\begin{equation}
		\vx_{t_{\ell - 1}} - \vmu = \vV\mat{1 - \ell T^{2}/(\lambda_{1}L^{2} + \ell^{2} T^{2}) & & \\ & \ddots & \\ & & 1 - \ell T^{2}/(\lambda_{D}L^{2} + \ell^{2}T^{2})}\vV^{\top}(\vx_{t_{\ell}} - \vmu).
	\end{equation}
	Notice that all entries in this diagonal matrix are in \((0, 1)\). Taking the \(\ell^{2}\)-norm and using a cheap yet tight upper-bound, it holds 
	\begin{equation}
		\norm{\vx_{t_{\ell - 1}} - \vmu}_{2} \leq \underbrace{\max_{i \in [D]}\bc{1 - \frac{\ell T^{2}}{\lambda_{i}L^{2} + \ell^{2}T^{2}}}}_{c_{\ell, L}} \cdot \norm{\vx_{t_{\ell}} - \vmu}_{2}.
	\end{equation}
	At first glance, this inequality seems to imply that the denoising iteration is \textit{very close}, conceptually, to a \textit{contraction mapping} such as power iteration (\eqref{eq:power_iteration}) sending the iterate \(\vx_{t}\) to the mean \(\vmu\), since \(c_{\ell, L}\) are all in \((0, 1)\) for any \(\ell\) and \(L\). However, this is false. To see why, observe that
	\begin{equation}
		\norm{\vx_{0} - \vmu}_{2} \leq \prod_{\ell = 1}^{L}c_{\ell, L} \cdot \norm{\vx_{T} - \vmu}_{2}.
	\end{equation}
	If each denoising update had similar properties to a contraction mapping, then \(\prod_{\ell = 1}^{L}c_{\ell, L}\) would go to \(0\) as  \(L \to \infty\), and the iterate would converge to \(\vmu\) extremely fast, rendering the sampler completely useless or vacuous. But in reality the denoising iteration is \textit{actually not a contraction mapping} since \(c_{\ell, L}\) changes with the index \(\ell\), so that \(\prod_{\ell = 1}^{L}c_{\ell, L}\) does not go to \(0\) as \(L \to \infty\), as demonstrated in \Cref{fig:not_contraction_mapping} (it is also possible to compute the limit analytically). In reality, the sampler ensures that the terminal iterates spread out amongst the support of the data distribution, enabling actual sampling from the data distribution, which is quite unlike the objective of power iteration (for example).
	\begin{figure}
		\centering 
		\includegraphics[width=0.75\textwidth]{\toplevelprefix/chapters/chapter3/figs/not_contraction_mapping.png}
		\caption{\small\textbf{The denoising iteration from \Cref{ex:single_gaussian_denoising} is not a contraction mapping.} Plot of the contraction coefficient products \(\prod_{\ell = 1}^{L}c_{\ell, L}\) in the case where the data are in \(D = 1\) , the time horizon \(T = 1\), and the data variance \(\lambda_{1} = a\), for varying values of \(a\). If the denoising iteration were (conceptually similar to) a contraction mapping, these products would go to \(0\), but instead they converge to non-zero values.}
		\label{fig:not_contraction_mapping}
	\end{figure}
\end{example}

Recall that we wanted to build the iterative denoising process to reduce the entropy of the iterates. While we did do this in a roundabout way by inverting a process which adds entropy, it is now time to pay the piper and confirm that our iterative denoising process reduces the entropy.

\begin{theorem}[Simplified Version of \Cref{thm:conditioning_reduces_entropy}]
	Suppose that \((\vx_{t})_{t \in [0, T]}\) obeys \eqref{eq:additive_gaussian_noise_model}. Then, under certain technical conditions on \(\vx\), for every \(s < t\) with \(s, t \in (0, T]\),
	\begin{equation}
		h(\Ex[\vx_{s} \mid \vx_{t}]) < h(\vx_{t}).
	\end{equation}
\end{theorem}
The full statement of the theorem, and the proof itself, requires some technicality, so it is postponed to \Cref{sub:denoising_entropy_decreases}.

\subsubsection{Computing the Denoiser}
The last thing we discuss here is that many times, we will \textit{not be able to compute} \(\bar{\vx}^{\ast}(t, \cdot)\) for any \(t\), since we do not have the distribution \(p_{t}\). But we can try to \textit{learn one from data}. Recall that the denoiser \(\bar{\vx}^{\ast}\) is defined in \eqref{eq:denoising_loss} as minimizing the mean-squared error \(\Ex\norm{\bar{\vx}(t, \vx_{t}) - \vx}_{2}^{2}\). We can use this mean-squared error as a loss or objective function to learn the denoiser. For example, we can parameterize \(\bar{\vx}(t, \cdot)\) by a neural network, writing it as \(\bar{\vx}_{\theta}(t, \cdot)\), and optimize the loss over the parameter space \(\Theta\):
\begin{equation}
	\min_{\theta \in \Theta}\Ex_{\vx, \vx_{t}}\norm{\bar{\vx}_{\theta}(t, \vx_{t}) - \vx}_{2}^{2}.
\end{equation}
The solution to this optimization problem, implemented via gradient descent or a similar algorithm, will give us a \(\bar{\vx}_{\theta^{\ast}}(t, \cdot)\) which is a good approximation to \(\bar{\vx}^{\ast}(t, \cdot)\) (at least if the training works) and which we will use as our denoiser.

What is a good architecture for this neural network \(\bar{\vx}_{\theta^{\ast}}(t, \cdot)\)? To answer this question, we will first attempt to understand the case of data \(\vx\) which are \textit{low-rank}, i.e., having low-rank support. The most canonical way to represent low-rank data is through a \textit{low-rank Gaussian}, and so we start our exploration by assuming that \(\vx\) is drawn (approximately) from a low-rank Gaussian:
\begin{equation}\label{eq:x_gaussian_distribution}
	\vx \sim \dNorm(\vzero, \vU\vU^{\top}),
\end{equation}
where \(\vU \in \O(D, P)\) is a tall orthogonal matrix with \(P \leq D\) (using zero-mean for essentially notational convenience).  The optimal denoiser under \eqref{eq:x_gaussian_distribution} is (from \Cref{example:denoising_gaussian_mixture}):
\begin{equation}\label{eq:optimal denoiser}
	\bar{\vx}^{\ast}(t, \vx_{t}) = \vU\vU^{\top}(\vU\vU^{\top} + t^{2}\vI)^{-1}\vx_{t}.
\end{equation}
The matrix  \(\vU\vU^{\top} + t^{2}\vI\) is a low-rank perturbation of the full-rank matrix \(t^{2}\vI\), and thus ripe for simplification via the \textit{Sherman-Morrison-Woodbury identity}, i.e., for matrices \(\vA, \vC, \vU, \vV\) such that \(\vA\) and \(\vC\) are invertible,
\begin{equation}\label{eq:sherman_morrison_woodbury}
	(\vA + \vU\vC\vV)^{-1} = \vA^{-1} - \vA^{-1}\vU(\vC^{-1} + \vV\vA^{-1}\vU)^{-1}\vV\vA^{-1}.
\end{equation}
We prove this identity in \Cref{exercise:sherman_morrison_woodbury_identity}. For now we apply this identity with \(\vA = t^{2}\vI\), \(\vU = \vU\), \(\vV = \vU^{\top}\), and \(\vC = \vI\), obtaining
\begin{align}
	(\vU\vU^{\top} + t^{2}\vI)^{-1} 
	&= \frac{1}{t^{2}}\vI - \frac{1}{t^{4}}\vU\bp{\vI + \frac{1}{t^{2}}\vU^{\top}\vU}^{-1}\vU^{\top} \\
	&= \frac{1}{t^{2}}\vI - \frac{1}{t^{4}\bp{1 + \frac{1}{t^{2}}}}\vU\vU^{\top} \\
	&= \frac{1}{t^{2}}\bp{\vI - \frac{1}{1 + t^{2}}\vU\vU^{\top}}.\label{eq:gaussian_lowrank_inverse_perturbation_simplified}
\end{align}
This implies
\begin{align}
	\vU\vU^{\top}(\vU\vU^{\top} + t^{2}\vI)^{-1}
	&= \frac{1}{t^{2}}\vU\vU^{\top}\bp{\vI - \frac{1}{1 + t^{2}}\vU\vU^{\top}} \\
	&= \frac{1}{t^{2}}\bp{1 - \frac{1}{1 + t^{2}}}\vU\vU^{\top} \\
	&= \frac{1}{1 + t^{2}}\vU\vU^{\top}.\label{eq:gaussian_lowrank_denoiser_simplified}
\end{align}
Our optimal denoiser is a (rescaled) linear projection onto the low-rank subspace supporting our data distribution. This is essentially the same thing as probabilistic PCA from \Cref{ch:classic} (cf \Cref{subsec:probabilistic PCA}). We can formalize this intuition by considering a class of subspace-parameterized denoisers:
\begin{equation}\label{eq:optimal nn}
	\bar{\vx}_{\vV}(t, \vx_{t}) = \frac{1}{1 + t^{2}}\vV\vV^{\top}\vx_{t},
\end{equation}
and showing that the denoising problem \(\min_{\vV \in \O(D, P)}\Ex_{\vx, \vx_{t}}\norm{\bar{\vx}_{\vV}(t, \vx_{t}) - \vx}_{2}^{2}\) obtains the same solution as probabilistic PCA (\Cref{subsec:probabilistic PCA}). Indeed, plugging our subspace denoiser into the training loss \eqref{eq:denoising_loss}, we obtain
\begin{equation}
	\min_{\vV \in \O(D, P)}\Ex_{\vx, \vx_{t}}\norm{\bar{\vx}_{\vV}(t, \vx_{t}) - \vx}_{2}^{2} = \Ex_{\vx, \vg}\norm*{\frac{1}{1 + t^{2}}\vV\vV^{\top}(\vx + t\vg) - \vx}_{2}^{2},
\end{equation}
where the equality is due to \eqref{eq:additive_gaussian_noise_model}. Conditioned on \(\vx\), we compute 
\begin{align}
  & \Ex_{\vg}\norm*{\frac{1}{1 + t^{2}}\vV\vV^{\top}(\vx + t\vg) - \vx}_{2}^{2} \\
  &= \norm*{\frac{1}{1 + t^{2}}\vV\vV^{\top}\vx - \vx}_{2}^{2} - \frac{t}{1+t^2}\Ex_{\vg} \ip*{\vx - \frac{1}{1 + t^{2}}\vV\vV^{\top}\vx ,  \vV\vV^{\top}\vg} \notag \\
  &\quad + \frac{t^2}{(1+t^2)^2} \Ex_{\vg}\norm*{\vV\vV^{\top} \vg}_{2}^{2} \\
  &= \norm*{\frac{1}{1 + t^{2}}\vV\vV^{\top}\vx - \vx}_{2}^{2} + \frac{t^2P}{(1+t^2)^2} 
\end{align}
where the second equality follows from \(\vg \sim \dNorm(\vzero, \vI)\) and \( \Ex_{\vg}\norm{\vV\vV^\top \vg}_{2}^{2} = \Ex_{\vg}\norm{\vV^\top \vg}_{2}^{2} = P\), since \(\vV \in \O(D, P)\). Therefore, the denoising problem is equivalent to 
\begin{align}
	\min_{\vV \in \O(D, P)} & \Ex_{\vx} \norm*{\frac{1}{1 + t^{2}}\vV\vV^{\top}\vx - \vx}_{2}^{2} \\ 
	& \qquad = \Ex_{\vx}\norm{\vx}_{2}^{2} + \bp{\frac{1}{(1 + t^{2})^{2}} - \frac{2}{1 + t^{2}}}\Ex_{\vx}\norm{\vV^{\top}\vx}_{2}^{2},
\end{align}
which in turn is equivalent to the problem
\begin{equation}\label{eq:PCA}
    \max_{\vV \in \O(D, P)}\Ex_{\vx}\norm{\vV^{\top}\vx}_{2}^{2}
\end{equation}
which was shown in \Cref{sub:pca} to be equivalent to the problem 
\begin{equation}
    \min_{\vV \in \O(D, P)}\Ex_{\vx}\norm{\vx - \vV\vV^{\top}\vx}_{2}^{2},
\end{equation}
i.e., the probabilistic PCA formulation \eqref{eq:PPCA}.

Thus, if we were to assume our data were drawn from a low-rank Gaussian, our denoising algorithm essentially would consist of many PCA iterations in a row. While this method sounds very simple, it can actually produce complex and relatively powerful denoising behaviors when applied to practical data, as demonstrated by PCANet \parencite{chan2015pcanet}. Nevertheless, in order to obtain a more general denoising operator that may work on complex and multi-modal data distributions, we will consider a more broad class of data distributions: the ubiquitous class of \textit{low-rank Gaussian mixture models},  whose denoiser we computed in \Cref{ex:gaussian_mixture} and \Cref{example:denoising_gaussian_mixture}. In the context of denoising, since Gaussian mixture models can approximate a very broad class of distributions arbitrarily well, in practice optimizing among the class of denoisers for Gaussian mixture models can potentially give us something close to the optimal denoiser for the real data distribution. 

Formally, we now assume that \(\vx\) is drawn from a low-rank Gaussian mixture model \(\frac{1}{K}\sum_{k = 1}^{K}\dNorm(\vzero, \vU_{k}\vU_{k}^{\top})\), where  \(\vU_{k} \in \O(D, P)\) are tall orthogonal matrices whose columns span the \(P\)-dimensional subspaces which each support one of the \(K\) low-rank Gaussian components. We write
\begin{equation}\label{eq:MoG1}
	\vx \sim \frac{1}{K}\sum_{k = 1}^{K}\dNorm(\vzero, \vU_{k}\vU_{k}^{\top}).
\end{equation}
Then the optimal denoiser under \eqref{eq:additive_gaussian_noise_model} is (from \Cref{example:denoising_gaussian_mixture})
\begin{equation} 
	\bar{\vx}^{\ast}(t, \vx_{t}) = \sum_{k = 1}^{K}\frac{\phi(\vx_{t} ; \vzero,
	\vU_{k}\vU_{k}^{\top} + t^{2}\vI)}{\sum_{i = 1}^{K}\phi(\vx_{t} ; \vzero, \vU_{i}\vU_{i}^{\top} + t^{2}\vI)}\cdot\bp{\vU_{k}\vU_{k}^{\top}(\vU_{k}\vU_{k}^{\top} + t^{2}\vI)^{-1}\vx_{t}}.
\end{equation}
Applying the same Sherman-Morrison-Woodbury identity as in \eqref{eq:gaussian_lowrank_inverse_perturbation_simplified}, we obtain
\begin{equation}
	(\vU_{k}\vU_{k}^{\top} + t^{2}\vI)^{-1} = \frac{1}{t^{2}}\bp{\vI - \frac{1}{1 + t^{2}}\vU_{k}\vU_{k}^{\top}}.
\end{equation}
Then we can compute the posterior probabilities as follows. Note that since \(\vU_{k}\)'s are all orthogonal, \(\det(\vU_{k}\vU_{k}^{\top} + t^{2}\vI)\) are all the same for each \(k\). So
\begin{align}
	\frac{\phi(\vx_{t} ; \vzero, \vU_{k}\vU_{k}^{\top} + t^{2}\vI)}{\sum_{i
	= 1}^{K}\phi(\vx_{t} ; \vzero, \vU_{i}\vU_{i}^{\top} + t^{2}\vI)} 
	&= \frac{\exp\rp{-\frac{1}{2}\vx_{t}^{\top}(\vU_{k}\vU_{k}^{\top} + t^{2}\vI)^{-1}\vx_{t}}}{\sum_{i = 1}^{K}\exp\rp{-\frac{1}{2}\vx_{t}^{\top}(\vU_{i}\vU_{i}^{\top} + t^{2}\vI)^{-1}\vx_{t}}} \\
	&= \frac{\exp\rp{-\frac{1}{2t^{2}}\vx_{t}^{\top}\bp{\vI - \frac{1}{1 + t^{2}}\vU_{k}\vU_{k}^{\top}}\vx_{t}}}{\sum_{i = 1}^{K}\exp\rp{-\frac{1}{2t^{2}}\vx_{t}^{\top}\bp{\vI - \frac{1}{1 + t^{2}}\vU_{i}\vU_{i}^{\top}}\vx_{t}}} \\
	&= \frac{\exp\rp{-\frac{1}{2t^{2}}\norm{\vx_{t}}_{2}^{2} + \frac{1}{2t^{2}(1 + t^{2})}\norm{\vU_{k}^{\top}\vx_{t}}_{2}^{2}}}{\sum_{i = 1}^{K}\exp\rp{-\frac{1}{2t^{2}}\norm{\vx_{t}}_{2}^{2} + \frac{1}{2t^{2}(1 + t^{2})}\norm{\vU_{i}^{\top}\vx_{t}}_{2}^{2}}} \\
	&= \frac{\exp\rp{-\frac{1}{2t^{2}}\norm{\vx_{t}}_{2}^{2}}\exp\rp{\frac{1}{2t^{2}(1 + t^{2})}\norm{\vU_{k}^{\top}\vx_{t}}_{2}^{2}}}{\exp\rp{-\frac{1}{2t^{2}}\norm{\vx_{t}}_{2}^{2}}\sum_{i = 1}^{K}\exp\rp{\frac{1}{2t^{2}(1 + t^{2})}\norm{\vU_{i}^{\top}\vx_{t}}_{2}^{2}}} \\
	&= \frac{\exp\rp{\frac{1}{2t^{2}(1 + t^{2})}\norm{\vU_{k}^{\top}\vx_{t}}_{2}^{2}}}{\sum_{i = 1}^{K}\exp\rp{\frac{1}{2t^{2}(1 + t^{2})}\norm{\vU_{i}^{\top}\vx_{t}}_{2}^{2}}}.
\end{align}
This is a softmax operation weighted by the projection of \(\vx_{t}\) onto each subspace measured by \(\norm{\vU_{i}^{\top}\vx_{t}}_{2}\) (tempered by the temperature \(2t^{2}(1 + t^{2})\)). Meanwhile, the component denoisers can be written in the same way as \eqref{eq:gaussian_lowrank_denoiser_simplified} to obtain 
\begin{equation}
	\vU_{k}\vU_{k}^{\top}(\vU_{k}\vU_{k}^{\top} + t^{2}\vI)^{-1}\vx_{t} = \frac{1}{1 + t^{2}}\vU_{k}\vU_{k}^{\top}\vx_{t}.
\end{equation}
Putting these together, we have 
\begin{equation}\label{eq:gmm_lowrank_denoiser}
	\bar{\vx}^{\ast}(t, \vx_{t}) = \frac{1}{1 + t^{2}}\sum_{k = 1}^{K}\frac{\exp\rp{\frac{1}{2t^{2}(1 + t^{2})}\norm{\vU_{k}^{\top}\vx_{t}}_{2}^{2}}}{\sum_{i = 1}^{K}\exp\rp{\frac{1}{2t^{2}(1 + t^{2})}\norm{\vU_{i}^{\top}\vx_{t}}_{2}^{2}}}\vU_{k}\vU_{k}^{\top}\vx_{t},
\end{equation}
i.e., a projection of \(\vx_{t}\) onto each of \(K\) subspaces, weighted by a soft-max operation of a quadratic function of \(\vx_{t}\). This functional form is similar to an \textit{attention mechanism} in a transformer architecture! As we will see in \Cref{ch:representation}, this is no coincidence at all; the deep link between denoising and lossy compression (to be covered in \Cref{sec:lossy_compression}) makes transformer denoisers so effective in practice. And so overall, our Gaussian mixture model theory motivates the use of transformer-like neural networks for denoising.

Overall, the learned denoiser architecture corresponds to an (implicit parametric) encoding scheme of the given data, since it can be used to denoise/project onto the data distribution. Training a denoiser is equivalent to finding a better coding scheme, and this partially fulfills one of the desiderata (i.e., pursuing a distribution by searching for schemes with better coding rates) at the end of \Cref{sub:min_entropy}. In the sequel, we will discuss how to fulfill the other.



\subsection{Sampling a Distribution via Iterative Denoising}\label{sub:sampling_denoising}

Remember that at the end of \Cref{sub:min_entropy}, we discussed a pair of desiderata for pursuing a distribution with low-dimensional structure. The first such desideratum is to start with a normal distribution, say with high entropy, and gradually reduce its entropy until it reaches the distribution of the data. We will call this procedure \textit{sampling} since we are generating new samples. It is now time for us to discuss how to do this with the toolkit we have built up.

We know how to denoise very noisy samples \(\vx_{T}\) to attain approximations \(\hat{\vx}\) that have similar distributions to the target random variable \(\vx\). But the desideratum says that, to sample, we want to start with a template distribution with \textit{no} influence from the distribution of \(\vx\) and use the denoiser to guide the iterates towards the distribution of \(\vx\). How can we do this? One way is motivated as follows:
\begin{equation}\label{eq:ve_converge_to_large_gaussian}
	\frac{\vx_{T}}{T} = \frac{\vx + T\vg}{T} = \frac{\vx}{T} + \vg \xrightarrow{T\to +\infty} \vg \sim \dNorm(\vzero, \vI).
\end{equation}
Thus, \(\vx_{T} \approx \dNorm(\vzero, T^{2}\vI)\). This approximation is quite good for almost all practical distributions, and visualized in \Cref{fig:xT_vs_noise}.
\begin{figure}
	\centering 
	\includegraphics[width=0.75\textwidth]{\toplevelprefix/chapters/chapter3/figs/xT_vs_noise.png}
	\caption{\small\textbf{Visualizing \(\vx_{T}\) versus \(\dNorm(\vzero, T^{2}\vI)\).} \textit{Left:} A plot of Gaussian mixture model data \(\vx\). \textit{Right:} A plot of \(\vx\) as well as \(\vx_{T}\) and an independent sample of \(\dNorm(\vzero, T^{2}\vI)\), for \(T = 10\). On the right plot, \(\vx\) is plotted in the same colors as the left: however, samples from \(\vx_{T}\) and \(\dNorm(\vzero, T^{2}\vI)\) are both much larger, on average, than samples from \(\vx\), and so it appears much smaller because of the scaling. Despite this large blow-up, we clearly observe the similarities in the distributions of \(\vx_{T}\) and \(\dNorm(\vzero, T^{2}\vI)\).}
	\label{fig:xT_vs_noise}
\end{figure}

So, discretizing \([0, T]\) into \(0 = t_{0} < t_{1} < \cdots < t_{L} = T\) uniformly using \(t_{\ell} = T\ell / L\) (as in the previous section), one possible way to sample from pure noise is:
\begin{itemize}
	\item Sample \(\hat{\vx}_{T} \sim \dNorm(\vzero, T^{2}\vI)\) (i.i.d.~of everything else) 
	\item Run the denoising iteration as in \Cref{sub:intro_diffusion_denoising}, i.e.,
		\begin{equation}\label{eq:denoising-iteration-basic}
		\hat{\vx}_{t_{\ell - 1}} = \bp{1 - \frac{1}{\ell}}\cdot\hat{\vx}_{t_{\ell}} + \frac{1}{\ell}\cdot\bar{\vx}^{\ast}(t_{\ell}, \hat{\vx}_{t_{\ell}}).
	\end{equation}
	\item Output \(\hat{\vx} = \hat{\vx}_{0}\).
\end{itemize}
This conceptually is all there is behind \textit{diffusion models}, which transform noise into data samples in accordance with the first desideratum. However, there are a few steps left to take before we get models which can actually sample from real data distributions like images given practical resource constraints. In the sequel, we will introduce and motivate several such steps.

\paragraph{Step 1: different discretizations.} The first step we do is motivated by the following point: \textit{we do not need to spend so many denoising iterations} at large  \(t\). If we look at \Cref{fig:ve_gmm_denoising}, we observe that the first \(200\) or \(300\) iterations out of the \(500\) iterations of the sampling process are just spent contracting the noise towards the data distribution as a whole, before the remaining iterations push the samples towards a subspace. Given a fixed iteration count \(L\), this signals that we should spend more timesteps \(t_{\ell}\) near \(t = 0\) compared to \(t = T\). During sampling (and training), we can therefore use another discretization of \([0, T]\) into \(0 \leq t_{0} < t_{1} < \cdots < t_{L} \leq T\), such as an \textit{exponential discretization}:
\begin{equation}\label{eq:denoising_exponential_discretization}
	t_{\ell} = C_{1}(e^{C_{2}\ell} - 1), \qquad \forall \ell \in \{0, 1, \dots, L\}
\end{equation}
where \(C_{1}, C_{2} > 0\) are constants which can be tuned for optimal performance in practice; theoretical analysis will often specify such optimal constants as well. Then the denoising/sampling iteration becomes 
\begin{equation}\label{eq:denoising_iteration}
	\hat{\vx}_{t_{\ell - 1}} \doteq \frac{t_{\ell - 1}}{t_{\ell}}\hat{\vx}_{t_{\ell}} + \bp{1 - \frac{t_{\ell - 1}}{t_{\ell}}}\bar{\vx}^{\ast}(t_{\ell}, \hat{\vx}_{t_{\ell}}),
\end{equation}
with, again, \(\hat{\vx}_{t_{L}} \sim \dNorm(\vzero, t_{L}^{2}\vI)\).

\begin{figure}[t]
	\centering 
	\includegraphics[width=\textwidth]{\toplevelprefix/chapters/chapter3/figs/vp_gmm_denoising.png}
	\caption{\small\textbf{Denoising a mixture of Gaussians using the VP diffusion process.} We use the same figure setup and data distribution as \Cref{fig:ve_gmm_denoising}. Note that compared to \Cref{fig:ve_gmm_denoising}, the noise distribution is much more concentrated around the origin.}
	\label{fig:vp_gmm_denoising}
\end{figure}

\paragraph{Step 2: different noise models.} The second step is to consider slightly different models compared to \eqref{eq:additive_gaussian_noise_model}. The basic motivation for this is as follows. In practice, the noise distribution \(\dNorm(\vzero, t_{L}^{2}\vI)\) becomes an increasingly poor estimate of the true covariance in high dimensions, i.e., \eqref{eq:ve_converge_to_large_gaussian} becomes an increasingly worse approximation, especially with anisotropic high-dimensional data. The increased distance between \(\dNorm(\vzero, t_{L}^{2}\vI)\) and the true distribution of \(\vx_{t_{L}}\) may cause the denoiser to perform worse in such circumstances. Theoretically, \(\vx_{t_{L}}\) never converges to any distribution as \(t_{L}\) increases, so this setup is difficult to analyze end-to-end. In this case, our remedy is to \textit{simultaneously add noise and shrink the contribution of \(\vx\), such that \(\vx_{T}\) converges as \(T \to \infty\)}. The rate of added noise is denoted \(\sigma \colon [0, T] \to \R_{\geq 0}\), and the rate of shrinkage is denoted \(\alpha \colon [0, T] \to \R_{\geq 0}\), such that \(\sigma\) is \textit{increasing} and \(\alpha\) is (not strictly) \textit{decreasing}, and 
\begin{equation}\label{eq:gen_additive_gaussian_noise_model}
	\vx_{t} \doteq \alpha_{t}\vx + \sigma_{t}\vg, \qquad \forall t \in [0, T].
\end{equation}
The previous setup has \(\alpha_{t} = 1\) and \(\sigma_{t} = t\), and this is called the \textit{variance-exploding (VE) process}. A popular choice which decreases the contribution of \(\vx\), as we described originally, has \(T = 1\) (so that \(t \in [0, 1]\)), \(\alpha_{t} = \sqrt{1 - t^{2}}\) and \(\sigma_{t} = t\); this is the \textit{variance-preserving (VP) process}. Note that under the VP process, \(\vx_{1} \sim \dNorm(\vzero, \vI)\) exactly, so we can just sample from this standard distribution and iteratively denoise. As a result, the VP process is much easier to analyze theoretically and more stable empirically.\footnote{Why use the whole \(\alpha, \sigma\) setup? As we will see in \Cref{exercise:implement_denoising_processes}, it encapsulates and unifies many proposed processes, including the recently popular so-called \textit{flow matching} process. Despite this, the VE and VP processes are still the most popular empirically and theoretically (so far), and so we will consider them in this Section.} 



With this more general setup, Tweedie's formula \eqref{eq:tweedie} becomes 
\begin{equation}\label{eq:gen_tweedie}
	\Ex[\vx \mid \vx_{t}] = \frac{1}{\alpha_{t}}\bp{\vx_{t} + \sigma_{t}^{2}\nabla \log p_{t}(\vx)}.
\end{equation}
The denoising iteration \eqref{eq:denoising_iteration} becomes
\begin{equation}\label{eq:gen_denoising_iteration}
	\hat{\vx}_{t_{\ell - 1}} = \frac{\sigma_{t_{\ell - 1}}}{\sigma_{t_{\ell}}}\hat{\vx}_{t_{\ell}} + \bp{\alpha_{t_{\ell - 1}} - \frac{\sigma_{t_{\ell - 1}}}{\sigma_{t_{\ell}}}\alpha_{t_{\ell}}}\bar{\vx}^{\ast}(t_{\ell}, \hat{\vx}_{t_{\ell}}).
\end{equation}

\begin{example}[Gaussian Mixture Model] 
	For a Gaussian mixture model (see \Cref{ex:gaussian_mixture}), its denoiser \eqref{eq:gmm_bayes_optimal_denoiser} becomes 
	\begin{equation}\label{eq:gen_gmm_bayes_optimal_denoiser}
		\bar{\vx}^{\ast}(t, \vx_{t}) = \sum_{k = 1}^{K}\frac{\pi_{k}\phi(\vx_{t}
		; \alpha_{t}\vmu_{k}, \alpha_{t}^{2}\vSigma_{k}
		+ \sigma_{t}^{2}\vI)}{\sum_{i = 1}^{K}\pi_{i}\phi(\vx_{t} ; \alpha_{t}\vmu_{i}, \alpha_{t}^{2}\vSigma_{i} + \sigma_{t}^{2}\vI)}\cdot\bp{\vmu_{k} + \alpha_{t}\vSigma_{k}(\alpha_{t}^{2}\vSigma_{k} + \sigma_{t}^{2}\vI)^{-1}(\vx_{t} - \alpha_{t}\vmu_{k})}.
	\end{equation}
	\Cref{fig:vp_gmm_denoising} demonstrates iterations of the sampling procedure. 
\end{example}

Note that the denoising iteration \eqref{eq:gen_denoising_iteration} gives a sampling algorithm called the DDIM (``Denoising Diffusion Implicit Model'') sampler \cite{song2020denoising}, and is one of the most popular sampling algorithms used today in diffusion models. We summarize it here in \Cref{alg:iterative_denoising}.



\begin{algorithm}
	\caption{Sampling using a denoiser.}
	\label{alg:iterative_denoising}
	\begin{algorithmic}[1]
		\Require{An ordered list of timesteps \(0 \leq t_{0} < \cdots < t_{L} \leq T\) to use for sampling.}
		\Require{A denoiser \(\bar{\vx} \colon \{t_{\ell}\}_{\ell = 1}^{L} \times \R^{D} \to \R^{D}\).}
		\Require{Scale and noise level functions \(\alpha, \sigma \colon \{t_{\ell}\}_{\ell = 0}^{L} \to \R_{\geq 0}\).}
		\Ensure{A sample \(\hat{\vx}\), approximately from the distribution of \(\vx\).}
		\Function{DDIMSampler}{$\bar{\vx}, (t_{\ell})_{\ell = 0}^{L}$}
		\State{Initialize \(\tilde{\vx}_{t_{L}} \sim\) approximate distribution of \(\vx_{t_{L}}\)} \Comment{VP \(\implies \dNorm(\vzero, \vI)\), VE \(\implies \dNorm(\vzero, t_{L}^{2}\vI)\).}
		\For{\(\ell = L, L - 1, \dots, 1\)}
		\State{Compute
			\begin{equation*}
				\hat{\vx}_{t_{\ell - 1}} \doteq \frac{\sigma_{t_{\ell - 1}}}{\sigma_{t_{\ell}}}\hat{\vx}_{t_{\ell}} + \bp{\alpha_{t_{\ell - 1}} - \frac{\sigma_{t_{\ell - 1}}}{\sigma_{t_{\ell}}}\alpha_{t_{\ell}}}\bar{\vx}(t_{\ell}, \hat{\vx}_{t_{\ell}})
			\end{equation*}
		}
		\EndFor
		\State{\Return{\(\hat{\vx}_{t_{0}}\)}}
		\EndFunction
	\end{algorithmic}
\end{algorithm}

\paragraph{Step 3: optimizing training pipelines.} If we use the procedure dictated by \Cref{sub:intro_diffusion_denoising} to learn a separate denoiser \(\bar{\vx}(t, \cdot)\) for each time \(t\) to be used in the sampling algorithm, \textit{we would have to learn \(L\) separate denoisers!} This is highly inefficient---the usual case is that we have to train \(L\) separate neural networks, taking up \(L\) times the training time and storage memory, and then be locked into using these timesteps for sampling forever. Instead, we can \textit{train a single neural network} to denoise across all times \(t\), taking as input the continuous variables \(\vx_{t}\) and \(t\) (instead of just \(\vx_{t}\) before). Mechanically, our training loss averages over \(t\), i.e., solves the following problem:
\begin{equation}\label{eq:training loss}
	\min_{\theta}\Ex_{t, \vx, \vx_{t}}\norm{\bar{\vx}_{\theta}(t, \vx_{t}) - \vx}_{2}^{2}.
\end{equation}
Similar to Step 1, where we used more timesteps closer to \(t = 0\) to ensure a better sampling process, we may want to ensure that the denoiser is higher quality closer to \(t = 0\), and thereby \textit{weight the loss} so that \(t\) near \(0\) has higher weight. Letting \(w_{t}\) be the weight at time \(t\), the weighted loss would look like
\begin{equation}\label{eq:parametric_denoising_loss}
	\min_{\theta}\Ex_{t}w_{t}\Ex_{\vx, \vx_{t}}\norm{\bar{\vx}_{\theta}(t, \vx_{t}) - \vx}_{2}^{2}.
\end{equation}
One reasonable choice of weight in practice is \(w_{t} = \alpha_{t}/\sigma_{t}\). The precise reason will be covered in the next paragraph, but generally it serves to up-weight the losses corresponding to \(t\) near \(0\) while still remaining reasonably numerically stable. Also, of course, we cannot compute the expectation in practice, so we use the most straightforward Monte-Carlo average to estimate it. The series of changes made here have several conceptual and computational benefits: we do not need to train multiple denoisers, we can train on one set of timesteps and sample using a subset (or others entirely), etc. The full pipeline is discussed in \Cref{alg:learning_denoiser}.

\begin{algorithm}
	\begin{algorithmic}[1]
		\Require{Dataset \(\cD \subseteq \R^{D}\).}
		\Require{An ordered list of timesteps \(0 \leq t_{0} < \cdots < t_{L} \leq T\) to use for sampling.}
		\Require{A weighting function \(w \colon \{t_{\ell}\}_{\ell = 1}^{L} \to \R_{\geq 0}\).}
		\Require{Scale and noise level functions \(\alpha, \sigma \colon \{t_{\ell}\}_{\ell = 0}^{L} \to \R_{\geq 0}\).}
		\Require{A parameter space \(\Theta\) and a denoiser architecture \(\bar{\vx}_{\theta}\).}
		\Require{An optimization algorithm for the parameters.}
		\Require{The number of optimization iterations \(M\).}
		\Require{The number of Monte-Carlo draws \(N\) per iteration (to approximate the expectation in \eqref{eq:parametric_denoising_loss})}
		\Ensure{A trained denoiser \(\bar{\vx}_{\theta^{\ast}}\).}

		\Function{TrainDenoiser}{$\cD, \Theta$}
		\State{Initialize \(\theta^{(1)} \in \Theta\)}
		\For{\(m \in [M]\)}
		\For{\(i \in [N]\)}
		\State{\(\vx^{(m, i)} \sim \cD\)} \Comment{Draw a sample from the dataset.}
		\State{\(t^{(m, i)} \simiid \dUnif(\{t_{\ell}\}_{\ell = 1}^{L})\)} \Comment{Sample a timestep.}
        \State{\(\alpha^{(m, i)} \doteq \alpha_{t^{(m, i)}}\)} \Comment{Compute the mean scale.}
        \State{\(\sigma^{(m, i)} \doteq \sigma_{t^{(m, i)}}\)} \Comment{Compute the noise scale.}
		\State{\(w^{(m, i)} \doteq w_{t^{(m, i)}}\)} \Comment{Compute the loss weight.}
		\State{\(\vg^{(m, i)} \simiid \dNorm(\vzero, \vI)\)} \Comment{Sample a noise vector.}
		\State{\(\vx_{t}^{(m, i)} \doteq \alpha^{(m, i)}\vx^{(m, i)} + \sigma^{(m, i)}\vg^{(m, i)}\)} \Comment{Compute the noised sample.}
		\EndFor
		\State{\(\hat{\cL}^{(m)} \doteq \displaystyle \frac{1}{N}\sum_{i = 1}^{N}w^{(m, i)}\norm{\vx^{(m, i)} - \bar{\vx}_{\theta^{(m)}}(t^{(m, i)}, \vx_{t}^{(m, i)})}_{2}^{2}\)} \Comment{Compute the loss.}
		\State{\(\theta^{(m + 1)} \doteq \texttt{OptimizationUpdate}^{(m)}(\theta^{(m)}, \nabla_{\theta^{(m)}}\hat{\cL}^{(m)})\)} \Comment{Update parameters.}
		\EndFor
		\State{\Return{\(\bar{\vx}_{\theta^{(M + 1)}}\)}}
		\EndFunction
	\end{algorithmic}
	\caption{Learning a denoiser from data.}
	\label{alg:learning_denoiser}
\end{algorithm}

\paragraph{(Optional) Step 4: changing the estimation target.} Note that it is common to instead reorient the whole denoising pipeline around \textit{noise predictors}, i.e., estimates of \(\Ex[\vg \mid \vx_{t}]\). In practice, noise predictors are slightly easier to train because their output is (almost) always of comparable size to a Gaussian random variable, so training is more numerically stable. Note that by \eqref{eq:gen_additive_gaussian_noise_model} we have 
\begin{equation}
	\vx_{t} = \alpha_{t}\Ex[\vx \mid \vx_{t}] + \sigma_{t}\Ex[\vg \mid \vx_{t}] \implies \Ex[\vg \mid \vx_{t}] = \frac{1}{\sigma_{t}}\bp{\vx_{t} - \alpha_{t}\Ex[\vx \mid \vx_{t}]},
\end{equation}
Therefore any predictor for \(\vx\) can be turned into a predictor for \(\vg\) using the above relation, i.e.,
\begin{equation}
	\bar{\vg}(t, \vx_{t}) = \frac{1}{\sigma_{t}}\vx_{t} - \frac{\alpha_{t}}{\sigma_{t}}\bar{\vx}(t, \vx_{t}),
\end{equation}
and vice-versa. Thus a good network for estimating \(\bar{\vg}\) is the same as a good network for estimating \(\bar{\vx}\) \textit{plus a residual connection} (as seen in, e.g., transformers). Their losses are also the same as the denoiser, up to the factor of \(\alpha_{t}/\sigma_{t}\), i.e.,
\begin{equation}
	\Ex_{t}w_{t}\Ex_{\vg, \vx_{t}}\norm{\vg - \bar{\vg}(t, \vx_{t})}_{2}^{2} = \Ex_{t}w_{t}\frac{\alpha_{t}^{2}}{\sigma_{t}^{2}}\Ex_{\vx, \vx_{t}}\norm{\vx - \bar{\vx}(t, \vx_{t})}_{2}^{2}.
\end{equation}
For the sake of completeness we will mention that other targets have been proposed for different tasks, e.g., \(\Ex[\odv*{\vx_{t}}{t} \mid \vx_{t}]\) (called \(v\)-prediction or velocity prediction), etc., but denoising and noise prediction remain commonly used. Throughout the rest of this book we will only consider denoising.

\subsubsection{Sampling Error Rate of the Denoising Process.}
We have made lots of changes to our original platonic noising/denoising process. To assure ourselves that the new process still works in practice, we can compute numerical examples (such as \Cref{fig:vp_gmm_denoising}). To assure ourselves that it is theoretically sound, we can prove a \textit{bound on the error rate} for the sampling algorithm, which shows that the error rate is small. We will now furnish such a rate from the literature, which shows that the output distribution of the sampler converges in the so-called \textit{total variation (TV) distance} to the true distribution. The TV distance is defined between two random variables \(\vx\) and \(\vy\) as:
\begin{equation}
	\TV(\vx, \vy) \doteq \sup_{A \subseteq \R^{d}}\abs*{\Pr[\vx \in A] - \Pr[\vy \in A]}.
\end{equation}
If \(\vx\) and \(\vy\) are very close (uniformly), then the supremum will be small. So the TV distance measures the closeness of random variables. (It is indeed a metric, as the name suggests; the proof is an exercise.)

\begin{theorem}[\citep{li2024d} Theorem 1, Simplified]\label{thm:diffusion_sampler_convergence}
	Suppose that \(\Ex\norm{\vx}_{2} < \infty\).  If \(\vx\) is denoised according to the VP process with an exponential discretization\footnote{The precise definition is rather lengthy in our notation and only defined up to various absolute constants, so we omit it here for brevity. Of course it is in the original paper \citep{li2024d}.} as in \eqref{eq:denoising_exponential_discretization}, the output \(\hat{\vx}\) of \Cref{alg:iterative_denoising} satisfies the total variation bound
	\begin{equation}\label{eq:diffusion_sampling_error}
		\TV(\vx, \hat{\vx}) = \tilde{\cO}\rp{\underbrace{\frac{D}{L}}_{\textbf{discretization error}} + \underbrace{\sqrt{\frac{1}{L}\sum_{\ell = 1}^{L}\frac{\alpha_{t_{\ell}}^{2}}{\sigma_{t_{\ell}}^{4}}\Ex_{\vx, \vx_{t_{\ell}}}\norm{\bar{\vx}^{\ast}(t_{\ell}, \vx_{t_{\ell}}) - \bar{\vx}(t_{\ell}, \vx_{t_{\ell}})}_{2}^{2}}}_{\textbf{average excess error of the denoiser}}}
	\end{equation}
	where \(\bar{\vx}^{\ast}\) is the Bayes optimal denoiser for \(\vx\), and \(\tilde{\cO}\) is a version of the big-\(\cO\) notation which ignores logarithmic factors in \(L\).
\end{theorem}
The very high-level proof technique is, as discussed earlier, to bound the error at each step, distinguish the error sources (between discretization and denoiser error), and carefully ensure that the errors do not accumulate too much (or even cancel out). 

Note that if \(L \to \infty\) and we correctly learn the Bayes optimal denoiser \(\bar{\vx} = \bar{\vx}^{\ast}\) (such that the excess error is \(0\)), then the sampling process in \Cref{alg:iterative_denoising} yields a \textit{perfect (in distribution) inverse} of the noising process, since the error rate in \Cref{thm:diffusion_sampler_convergence} goes to \(0\),\footnote{There are similar results for VE processes, though none are as sharp as this to our knowledge.} as heuristically argued previously.

\begin{remark}
	What if the data is low-dimensional, say supported on a \textit{low-rank subspace} of the high dimensional space \(\R^{D}\)? If the data distribution is compactly supported---say if the data is normalized to the unit hypercube, which is often ensured as a pre-processing step for real data such as images---it is possible to do better. Namely, the authors of \cite{li2024d} also define a measure of \textit{approximate intrinsic dimension} using the asymptotics of the so-called covering number, which is extremely similar in intuition (if not in implementation) to the rate distortion function presented in the next Section. Then they show that using a particular small modification of the DDIM sampler in \Cref{alg:iterative_denoising} (i.e., slightly perturbing the update coefficients), the discretization error becomes
	\begin{equation}
		\tilde{\cO}\rp{\frac{\text{approximate intrinsic dimension}}{L}}
	\end{equation}
	instead of \(\frac{D}{L}\) like it was in \Cref{thm:diffusion_sampler_convergence}. Therefore, using this modified algorithm, \(L\) does not have to be too large even as \(D\) reaches the thousands or millions, since real data have low-dimensional structure. However in practice we use the DDIM sampler instead, so \(L\) should have a mild dependence on \(D\) to achieve consistent error rates. The exact choice of \(L\) trades off between the computational complexity (e.g., runtime or memory consumption) of sampling and the statistical complexity of learning a denoiser for low-dimensional structures. The value of \(L\) is often different at training time (where a larger \(L\) allows better coverage of the interval \([0, T]\), which helps the network learn a relationship which generalizes over \(t\)) and sampling time (where \(L\) being smaller means more efficient sampling). One can even pick the timesteps adaptively at sampling time in order to optimize this tradeoffs \cite{bao2022analytic}.
\end{remark}

\begin{remark}
	Various other works define the reverse process as moving backward in the time index \(t\) using an explicit difference equation, or differential equation in the limit \(L \to \infty\), or forward in time using the transformation \(\vy_{t} = \vx_{T - t}\), such that if \(t\) increases then \(\vy_{t}\) becomes closer to \(\vx_{0}\). In this work we strive to keep consistency: we move forward in time to noise, and backward in time to denoise. If you are reading another work which is not clear on the time index, or trying to implement an algorithm which is similarly unclear, there is one way to do it right every time: the sampling process should always have a \textit{positive} coefficient on both the denoiser term and the current iterate when moving from step to step. But in general many papers define their own notation and it is not user-friendly.
\end{remark}

\begin{remark}
	The theory presented at the end of the last \Cref{sub:intro_diffusion_denoising} seems to suggest (loosely speaking) that in practice, using a transformer-like network is a good choice for learning or approximating a denoiser. This is reasonable, but what is the problem with using any old neural network (such as a multi-layer perceptron (MLP)) and just trying to scale it up to infinity? To observe the problem with this, let us look at another special case of the Gaussian mixture model studied in \Cref{example:denoising_gaussian_mixture}. Namely, the \textit{empirical distribution} on vectors \((\vx^{(i)})_{i = 1}^{N} \subseteq \R^{D}\) is an instance of a degenerate Gaussian mixture model, with \(K = N\) components \(\dNorm(\vx^{(i)}, \vzero)\) sampled with equal probability \(\pi_{i} = \frac{1}{N}\). In this case the Bayes optimal denoiser is
	\begin{equation}\label{eq:memorizing_denoiser}
		\bar{\vx}^{\star}(t, \vx_{t}) = \sum_{i = 1}^{N}\frac{\exp({-\|\vx_{t} - \alpha_{t}\vx^{(i)}\|_{2}^{2}/(2\sigma_{t}^{2})})}{\sum_{j = 1}^{N}\exp({-\|\vx_{t} - \alpha_{t}\vx^{(j)}\|_{2}^{2}/(2\sigma_{t}^{2})})}\vx^{(i)}.
	\end{equation}
	This is a convex combination of the data \(\vx^{(i)}\), and the coefficients get ``sharper'' (i.e., closer to \(0\) or \(1\)) as \(t \to 0\). Notice that this denoiser \textit{optimally solves} the denoising optimization problem \eqref{eq:parametric_denoising_loss} when we compute the loss based on drawing \(\vx\) uniformly at random from a fixed finite dataset \(\vX = (\vx^{(i)})_{i = 1}^{N}\), which is a realistic model for training on a finite dataset. Thus, if our network architecture \(\bar{\vx}_{\theta}\) is expressive enough such that optimal denoisers of the above form \eqref{eq:memorizing_denoiser} may be well-approximated, then the learned denoiser may do just that. Then, our iterative denoising \Cref{alg:iterative_denoising} will sample exactly from the empirical distribution, re-generating samples in the training data, as certified by \Cref{thm:diffusion_sampler_convergence}. This is a bad sampler, not really more interesting than a database of all samples, and so it is important to understand how to avoid this in practice. The key is to come up with a network architecture which can well-approximate the true denoiser (say corresponding to a low-rank distribution as in \eqref{eq:gmm_lowrank_denoiser}) but not the empirical Bayesian denoiser as in \eqref{eq:memorizing_denoiser}. Some work has explored this fine line and why modern diffusion models, which use transformer- and convolutional-based network architectures, can memorize and generalize in different regimes \citep{kamb2024analytic,niedoba2024towards}.

	At a high level, a denoiser which memorizes all the training points, as in \eqref{eq:memorizing_denoiser}, corresponds to a parametric model of the distribution which has minimal coding rate, and achieves this by just coding every sample separately. We will discuss this problem (and seeming paradox with our initial desiderata at the end of \Cref{sub:min_entropy}) from the perspective of information theory in the next section.
\end{remark}



\section{Memorization and Generalization in Diffusion Models}

%\yima{Qing and Druv, a new section about memorization and generalization probably should go here. It should elucidate the role of quantization error (small noise) introduced by the network capacity... denoising to death means memorizing individual samples, hence leads to the topic of next chapter. } 

So far, we have discussed how to learn low-dimensional distributions via denoising by optimizing the population loss \eqref{eq:training loss}. However, in practice, as we only have finite training samples $(\vx^{(i)})_{i=1}^N$, we would minimize an empirical surrogate of \eqref{eq:training loss}, that is
\begin{equation}\label{eq:training loss-emp}
	\min_{\bm \theta} \frac{1}{N}\Ex_{t, \bm g \sim \mathcal{N}(\bm 0, \bm I)}\norm{\bar{\vx}_{\bm \theta}(t, \vx^{(i)}_t) - \vx^{(i)}}_{2}^{2},
\end{equation}
where $ \vx^{(i)}_t = \alpha_t \bm x^{(i)} + \sigma_t\bm g$ and $\bm g \sim \mathcal{N}(\bm 0, \bm I)$ as defined in \eqref{eq:gen_additive_gaussian_noise_model}.
Ideally, if our model $\bar{\vx}_{\bm \theta}$ has infinite model capacity, then it would merely memorize the training data itself $(\vx^{(i)})_{i=1}^N$ and thus learn the empirical distribution 
\begin{align*}
    p_{\text{emp}}(\vx) = \sum_{i=1}^N \delta(\vx -  \vx^{(i)}).
\end{align*}
Therefore, it raises a fundamental question regarding memorization vs. generalization. Specifically, \emph{how many training samples are required to learn the true data distribution $p_{\text{data}}$ instead of memorization?}

%sample complexity for learning the data distribution $p_{\text{data}}$. In our previous discussion, $p_{\text{data}}$ could be the mixture of low-rank Gaussians. 

 %While an overly expressive network can reproduce the empirical distribution exactly, practical diffusion models often capture the underlying structure of the data, enabling the generation of novel samples. 

%Building on the discussion of denoisers that simply memorize the training data, 

Based on the above discussion, we now turn to a deeper investigation of when diffusion models transition from memorization to generalization with respect to the number of training samples $N$.
Hypothetically, if a deep learning model $\bar{\vx}_{\bm \theta}$ has infinite model capacity, a perfectly learned neural network over \eqref{eq:training loss-emp} would simply learn $p_{\text{emp}}$ and memorize the training data without the ability to generate new samples. On the other hand, if we want $p_{\text{emp}}$ to approximate the true distribution $p_{\text{data}}$ to be $\epsilon$-close, in the worst case, the number of samples needed grows exponentially $O( (\frac{L}{\epsilon})^D )$ with respect to the data dimension $D$,\footnote{Here, $L$ denotes the Lipschitz constant of $p_{\text{data}}$.} implying that learning $p_{\text{data}}$ could potentially suffer from the curse of dimensionality.



However, in practice, real-world diffusion models, when trained on a finite training dataset $(\vx^{(i)})_{i=1}^N$ of natural images, can creatively generate new, sensible examples that differ significantly from the training data. In this section, we investigate the mechanisms and principles that enable the diffusion model to generalize when trained with finite samples and translate these insights into practical guidance for network design and training strategies.

% Despite the remarkable empirical success of diffusion models, their fundamental mechanisms are still poorly understood. In theory, the training of diffusion models is designed to learn the empirical distribution 
% \begin{align*}
%     p_{\text{emp}}(\vx) = \sum_{i=1}^N \delta(\vx - \vx^{(i)})
% \end{align*}
% over the training samples $\mathcal S= \{  \vx^{(i)} \}_{i=1}^N$ via neural networks (i.e., a summation of delta functions $\delta(\cdot)$ of the training data points), implying that a perfectly learned neural network would simply memorize the training data. But in practice, diffusion models can creatively generate new, sensible examples that differ significantly from the training data. This surprising discrepancy poses an important question about the generalizability of diffusion models.

\begin{figure}[t]
	\centering 
	\includegraphics[width=\textwidth]{\toplevelprefix/chapters/chapter3/figs/reproducibility.jpg}
	\caption{\small \textbf{Consistent model reproducibility across different diffusion models.} Visual content generated by three different diffusion models from the same initial noise (i.e., the noise in each row and column is the same) looks nearly identical. All models were trained on the CIFAR-10 dataset. 2a. Result of a denoising diffusion probabilistic model (DDPM). 2b. Result of consistency training (CT). 2c. Result of a vision transformer (ViT)-based architecture called U-ViT.
}
	\label{fig:vp_gmm_denoising}
\end{figure}

\subsubsection{Model Reproducibility Provides Insight into Generalization}  
Recent work \cite{pmlr-v235-zhang24cn} showed consistent model reproducibility among different diffusion models, which provides deep insight into the mystery of generalizability. In this intriguing, widespread phenomenon, different diffusion models produce remarkably similar outputs when initialized with the same noise and sampled via the same sampling procedure. This outcome holds across a variety of diffusion model frameworks, architectures, and training procedures (see \Cref{fig:vp_gmm_denoising}). Quantitatively, we can measure the reproducibility (RP) using the following probability metric:
\begin{align}
    \text{RP Score} := \mathbb P\left[\mathcal M_{\text{SSCD}}(\vx_1,\vx_2)>\tau\right],
\end{align}
where $\tau>0$ is some pre-defined threshold,\footnote{In \cite{pmlr-v235-zhang24cn}, the authors empirically set a large threshold $\tau =0.6$.} $\mathcal M_{\text{SSCD}}(\vx_1,\vx_2)$ measures the cosine similarity between $\vx_1$ and $\vx_2$ in the latent space of a self-supervised copy-detection (SSCD) neural descriptor, and $\vx_1$ and $\vx_2$ are generated by two different diffusion models starting from the same initial noise. Strong model reproducibility implies that different diffusion models learn the same score functions. However, we have not yet answered the critical question: Why are diffusion models able to generalize? To answer this question, we introduce a generalization (GL) score that quantifies the distinction between a generated sample $\vx$ and the training samples $(\vx^{(i)})_{i=1}^N$:
\begin{align}
    \text{GL Score} := 1- \mathbb P( \max_{i \in [N]} \mathcal M_{\text{SSCD}}(\vx,\vx^{(i)})>\tau).
\end{align}
Under the same setting, \Cref{fig:reproducibility-transition} depicts both the RP and GL scores for diffusion models that are trained with varying data sizes and model capacities, implying that these models learn two distinct distributions corresponding to two separate training regimes:
\begin{itemize}
    \item \textbf{Memorization regime:} When model capacity far exceeds training data size, all diffusion models overfit to the same empirical distribution $p_{\text{emp}}$ and memorize the training data, ultimately leading to strong model reproducibility without generalizability.
    \item \textbf{Generalization regime:} In contrast, when the model capacity is smaller than the size of the training data, diffusion models fail to memorize and then start to generate new samples. Here, reproducibility and generalizability co-emerge as all models learn the score function of the true distribution $p_{\text{data}}$ instead of memorizing the training data.
\end{itemize}
Moreover, within a fixed capacity, diffusion models transition from memorization to generalization as the number of training samples increases. While the seminal work \cite{pmlr-v235-zhang24cn} characterized this phase transition jointly using both the RP and GL scores, more recent work \cite{zhang2025understanding} unifies the metric by introducing the probability flow distance (PFD), which precisely characterizes the phase transition in a teacher-student setup. 

Empirically, this phase transition occurs with a finite number of samples (e.g., 10,000 samples from the CIFAR-10 dataset), without the curse of dimensionality. This is fundamentally different from the setting where the empirical distribution $p_{\text{emp}}$ interpolates the true data distribution $p_{\text{data}}$, in which case the required sample size can grow exponentially with the data dimension $D$.

\begin{figure}[t]
	\centering 
	\includegraphics[width=\textwidth]{\toplevelprefix/chapters/chapter3/figs/reproducibility-transition}
	\caption{\small \textbf{Co-emergence of reproducibility and generalizability} in the denoising diffusion probabilistic model, which is trained on a varying number of images from the CIFAR-10 dataset. 
}
	\label{fig:reproducibility-transition}
\end{figure}

\subsubsection{Towards a Mathematical Theory of Generalization}

To explain why diffusion models can generalize without suffering from the curse of dimensionality, recent work by Wang et al. \cite{wang2024diffusion} argues that they typically learn low-dimensional distributions rather than arbitrary ones. While real-world image datasets are high-dimensional in terms of raw pixel count and volume, extensive empirical evidence suggests they possess a significantly lower intrinsic dimension and reside on a union of low-dimensional manifolds \cite{brown2023verifying,gong2019intrinsic}. Specifically, Wang et al. \cite{wang2024diffusion} demonstrated that the required sample size for generalization scales linearly with this intrinsic dimension, even in high-dimensional ambient spaces.

%Diffusion models avoid the curse of dimensionality because they effectively learn low-dimensional distributions, where the required sample size for generalization scales linearly with the intrinsic dimension in high-dimensional space. While real-world image datasets are high-dimensional in terms of pixel count and data volume, extensive empirical evidence implies that they have a significantly lower intrinsic dimension and reside on the union of low-dimensional manifolds \cite{brown2023verifying,gong2019intrinsic}.


\begin{figure}[t]
\begin{center}
    \includegraphics[width = 1\linewidth]{\toplevelprefix/chapters/chapter3/figs/MoLRG_verify1.jpg}\vspace{-0.1in}
    \caption{\textbf{Comparison of images generated from the Gaussian, a mixture of low-rank Gaussians, and the distribution learned by diffusion models across different datasets.} Each row displays images generated from different distributions using the reverse-time ODE sampler, including the Gaussian, a mixture of low-rank Gaussians, and the distribution learned by U-Net. The columns represent images generated from the same initial noise. The results are shown for four datasets: FashionMNIST (top left ), MNIST (top right), CIFAR-10 (bottom left), and FFHQ (bottom right ).\label{fig:MoLRG_data_model_verfication}}%\qq{I feel we should remove the closest point in the training data, the 2nd row}  
\end{center}
\end{figure} 

Because the complexity of real-world image distributions makes it difficult to rigorously verify our claims on image datasets, Wang et al. \cite{wang2024diffusion} therefore consider a simplified setting of a mixture of low-rank Gaussians, as introduced in \eqref{eq:MoG1}, to mathematically characterize the generalization behavior.  As demonstrated in \Cref{fig:MoLRG_data_model_verfication}, the mixture of low-rank Gaussian distribution captures the essential low-dimensional structure of image distribution.

For simplicity, let us first introduce the result based a single low-rank Gaussian distribution, i.e., 
\begin{equation} 
	\vx \sim \dNorm(\vzero, \vU^*\vU^{*\top}),
\end{equation}
where $\bm U^* \in \O(D, P)$ denotes the ground-truth orthonormal basis. Accordingly, the training samples $( \vx^{(i)})_{i=1}^N$ are generated according to  
\begin{align}\label{eq:Ua+e} 
      \vx^{(i)} = \bm U^* \bm a_i + \bm e_i, 
\end{align}
where $\bm a_i \overset{i.i.d.}{\sim} \mathcal{N}(\bm 0, \bm I)$ is coefficients for each $i \in [N]$, and $\bm e_i \in \R^D$ is noise for each $i \in [N]$. Note that the noise $\bm e_i$ captures variability arising from real-world factors---such as measurement error, modeling mismatch, or deviations from an ideal low-rank structure. 

%According to \eqref{eq:training loss}, the empirical training loss of diffusion models is 
%\begin{equation}\label{eq:loss}
%	\min_{\theta} \frac{1}{N}\Ex_{t, \bm g \sim \mathcal{N}(\bm 0, \bm I)}\norm{\bar{\vx}_{\theta}(t, \vx^{(i)}_t) - \vx^{(i)}}_{2}^{2},
%\end{equation}
%where $ \vx^{(i)}_t = \alpha_t \bm x^{(i)} + \sigma_t\bm g$ and $\bm g \sim \mathcal{N}(\bm 0, \bm I)$ as defined in \eqref{eq:gen_additive_gaussian_noise_model}.  
Using \eqref{eq:gen_additive_gaussian_noise_model} and the same calculation in \eqref{eq:optimal denoiser}--\eqref{eq:gaussian_lowrank_denoiser_simplified}, we know that the optimal denoiser is  
\begin{align}
    \bar{\vx}^{\ast}(t, \vx_{t}) = \frac{\alpha_t}{\alpha_t^2 + \sigma_t^2}\bm U^*\bm U^{*T}\bm x_t.
\end{align}
Accordingly, we parameterize the network $\bar{\vx}_{\theta}$ as
\begin{align}\label{eqn:network-Gaussian}
    \bar{\vx}_{\theta}(t, \vx^{(i)}_t) = \frac{\alpha_t}{\alpha_t^2 + \sigma_t^2}\bm U\bm U^T\bm x_t,
\end{align}
where $\bm U \in \O(D, P)$ is learnable. By plugging \eqref{eqn:network-Gaussian} into \eqref{eq:training loss-emp} and using similar arguments in \eqref{eq:optimal nn}--\eqref{eq:PCA}, we can show that Problem \eqref{eq:training loss-emp} is equivalent to the classical PCA problem of the following form: 
    \begin{align}\label{eq:em loss}
        \max_{\bm U \in \O(D, P)} \sum_{i=1}^N \|\bm U^T \bm x^{(i)}\|^2. 
    \end{align}
This theorem shows that training diffusion models is equivalent to performing PCA on the training samples. Note that PCA is a classical and well-studied method for learning low-dimensional subspaces, whose optimal solution can be computed via singular value decomposition (SVD). This closed-form solution allows us to leverage existing results, such as Wedin's Theorem \cite{wedin1972perturbation}, to facilitate our analysis. Leveraging this equivalence and using the concentration inequalities, we can precisely determine the minimal number of samples required to learn a single low-rank Gaussian as follows: 

\begin{figure}[t]
\begin{center}
    \begin{subfigure}{0.43\textwidth}
        \includegraphics[width=1\textwidth]{\toplevelprefix/chapters/chapter3/figs//PCA.eps} 
    \caption{\bf\footnotesize PCA} 
    \end{subfigure}  \hfill
    \begin{subfigure}{0.43\textwidth}
    	\includegraphics[width = 1\linewidth]{\toplevelprefix/chapters/chapter3/figs//PCA_diffusion.eps} 
    \caption{\bf\footnotesize Diffusion Model} 
    \end{subfigure} \vspace{-0.1in}
    \caption{\textbf{Phase transition of learning a single low-rank Gaussian distribution with $K = 1$.} The $x$-axis is the number of training samples and $y$-axis is the dimension of subspaces. Darker pixels represent a lower empirical probability of success. We apply SVD and stochastic gradient descent to solve Problems \eqref{eq:PCA} and \eqref{eq:training loss-emp}, visualizing the results in (a) and (b), respectively.}  \label{fig:phase-transition-MoG-k=1} 
\end{center}
\end{figure}

\begin{theorem}\label{thm:2}
    Suppose that the training samples $\{\bm x^{(i)}\}_{i=1}^N$ are generated according to \eqref{eq:Ua+e}. Let $\hat{\bm U}$ denote an optimal solution of Problem (\ref{eq:training loss-emp}). The following statements hold: 
    \begin{itemize} 
        \item[i)] If $N \ge P$, it holds with probability at least $1 - 1/2^{N-P+1} - \exp\left(-c_2N\right) $ that any optimal solution $\hat{\bm U}$ satisfies
   \begin{align}\label{rst1:thm 2}
       \left\|\hat{\bm U}\hat{\bm U}^T - \bm U^*\bm U^{* T}\right\|_F \le \frac{c_1\sqrt{\sum_{i=1}^N \|\bm e_i\|^2}}{\sqrt{N} - \sqrt{P-1}},
   \end{align}
   where $c_1, c_2 > 0$ are constants. 
   \item[ii)] If $N < P$, there exists an optimal solution $\hat{\bm U} \in \O(D, P)$ such that with probability at least $1 - 1/2^{P-N+1} - \exp\left(-c_2^\prime P\right)$, 
   \begin{align}\label{rst2:thm 2}
       \left\|\hat{\bm U}\hat{\bm U}^T - \bm U^*\bm U^{* T}\right\|_F \ge \alpha - 
       \frac{c_1^\prime\sqrt{\sum_{i=1}^N \|\bm e_i\|^2}}{\sqrt{P} - \sqrt{N-1}},
   \end{align}
   where $\alpha := \sqrt{2\min\{P-N, D-P\}}$ and $c_1^\prime, c_2^\prime > 0$ are constants . 
    \end{itemize}
\end{theorem} 
This theorem demonstrates a phase transition from failure to success of diffusion models for learning the underlying distribution as the number of training samples increases. More precisely, when the number of training samples is larger than the intrinsic dimension of the subspace, i.e., $N \geq P$, any optimal solution $\hat{\bm U}$ recovers the basis of the underlying subspace with an approximation error depending on the noise level. Conversely, when $N < P$, optimizing the training loss fails to learn the underlying distribution. Please refer to \citep{wang2024diffusion} for the formal proof. 





Moreover, such insights can be generalized to learning a mixture of low-rank Gaussians \eqref{eq:MoG1}. Specifically, for simplicity, we consider a hard-max counterpart of \eqref{eq:gen_gmm_bayes_optimal_denoiser} for parameterizing the neural network $\bar{\vx}_{\bm \theta}$ as
    \begin{align}\label{eq:para MoG}
     \bar{\vx}_{\bm \theta}(\bm x_t, t)  = \dfrac{s_t}{s_t^2 + \gamma_t^2} \sum_{k=1}^K \hat{w}_k(\bm \theta; \bm x_0) \bm U_k\bm U_k^T\bm x_t, 
    \end{align}
where $\bm \theta = \{\bm U_k\}_{k=1}^K$, $\bm U = [\bm U_1,\dots,\bm U_K] \in \mathcal{O}^{D \times KP }$, and $\{ \hat{w}_k(\bm \theta; \bm x_0) \}_{k=1}^K$ are set as the following hard-max weights: 
\begin{align}\label{eq:wk1}
\hat{w}_k(\bm \theta; \bm x_0) = \begin{cases}
1,\ & \text{if}\ k=k_0, \\
0,\ & \text{otherwise},
\end{cases}
\end{align}
where $k_0 \in [K]$ is an index satisfying $\|\bm U_{k_0}^T\bm x_0\| \ge  \|\bm U_l^T\bm x_0\|$ for all $l \neq k_0 \in [K]$. 

\begin{theorem}\label{thm:3}
    Suppose that the neural network $\bar{\vx}_{\bm \theta}(\cdot,t)$ in Problem (\ref{eq:training loss-emp}) is parameterized into \eqref{eq:para MoG} for each $t \in [0,1]$, where $\hat{w}_k(\bm \theta, \bm x_0)$ is defined in \eqref{eq:wk1} for each $k \in [K]$. Then Problem (\ref{eq:training loss-emp}) is equivalent to the following subspace clustering problem: 
    \begin{align}\label{eq:SC}
        \max_{\bm \theta} \frac{1}{N} \sum_{k=1}^K \sum_{ i \in C_k(\bm \theta)} \|\bm U_k^T\bm x^{(i)}\|^2\qquad \mathrm{s.t.}\ \left[
            \bm U_1, \dots, \bm U_K \right] \in \mathcal{O}^{D \times KP}, 
    \end{align}
    where $C_k(\bm \theta) := \left\{ i\in [N]: \|\bm U_k^T\bm x^{(i)}\| \ge \|\bm U_l^T\bm x^{(i)}\|,\ \forall l \neq k \right\}$ for each $k \in [K]$.
\end{theorem} 
In this theorem, the constraint set $C_k(\bm \theta)$ ensures that each data point is assigned to the correct subspace—that is, the one onto which the norm of its projection is largest. Problem (\ref{eq:SC}) seeks to maximize the sum of squared norms of the projections of data points onto their respective assigned subspaces. With the network parameterization in \eqref{eq:para MoG}, \Cref{thm:3} shows that optimizing the training loss of diffusion models is equivalent to solving the subspace clustering problem \citep{vidal2011subspace,wang2022convergence}.  Notably, subspace clustering is a fundamental problem in unsupervised learning that aims to identify and group data points that lie in a union of low-dimensional subspaces in a high-dimensional space \citep{vidal2011subspace,lerman2018overview}. By showing an equivalence between training diffusion models and subspace clustering in \Cref{thm:3}, similar to \Cref{thm:2}, we can show that the minimum number of samples required for learning the underlying mixture of Gaussians depends linearly on the intrinsic subspace dimension for each Gaussian model. Similar to $K=1$ case, we can show the phase transition in \Cref{fig:phase-transition-MoG-k=2} that demonstrates our claim. For more details of the result, we refer readers to \cite{wang2024diffusion}.

\begin{figure}[t]
\begin{center}
    \begin{subfigure}{0.43\textwidth}
    	\includegraphics[width=1\textwidth]{\toplevelprefix/chapters/chapter3/figs/SC.eps} 
    \caption{\bf\footnotesize Subspace Clustering} 
    \end{subfigure}  
    \begin{subfigure}{0.43\textwidth}
    	\includegraphics[width = 1\linewidth]{\toplevelprefix/chapters/chapter3/figs/SC_diffusion.eps} 
    \caption{\bf\footnotesize Diffusion Model} 
    \end{subfigure}
    \caption{\textbf{Phase transition of learning a mixture of low-rank Gaussians with $K = 2$.} The $x$-axis is the number of training samples and $y$-axis is the dimension of subspaces. Darker pixels represent a lower empirical probability of success. We apply a subspace clustering method and stochastic gradient descent to solve Problems \eqref{eq:SC} and \eqref{eq:training loss-emp}, visualizing the results in (a) and (b), respectively. }  \label{fig:phase-transition-MoG-k=2} 
\end{center}
\end{figure}


\begin{remark}
Finally, we want to mention that the theoretical results we presented here are based on linear models. The simplified linear models provide the right intuition for characterizing the learning data distribution $p_{\text{data}}$, but they cannot fully characterize memorization. To deal with this, more recent work \cite{zhang2025generalization} considered the two-layer ReLU denoising autoencoder (DAE), which provides a unified understanding of memorization and generalization. Under the mixture of low-rank Gaussian distributions, the work characterized the global optimal solutions for training the two-layer ReLU network, showing that (i) memorization corresponds to the model storing raw training samples in the learned weights for encoding and decoding, whereas (ii) generalization arises when the model captures local data statistics. For more details of the result, we refer readers to \cite{zhang2025generalization}.    
\end{remark}


% According to \eqref{eq:optimal nn}-\eqref{eq:PCA}, seeking the optimal denoiser under the parameterization 


% To mathematically understand generalization behavior, let us consider a simple example of the mixture of low-rank Gaussian (MoLRG) distribution: 
% \begin{align}
%     \vx \sim \sum_{k=1}^K \pi_k \mathcal N(\mathbf 0, \bm U_k\bm U_k^\top),
% \end{align}
% where $\pi_k$ is the mixing proportion for each $k$ that satisfies $\sum_{k=1}^K \pi_k = 1$ with $\pi_k \geq 0$, and $\bm U_k$ is the orthonormal basis of the $k$th component for each $k$. Note that MoLRG data lie on a union of linear subspaces and locally approximate the low-dimensional manifolds of image data. Furthermore, the MoLRG model is amenable to mathematical analysis with closed-form score functions while also preserving the low-dimensional characteristics of image distributions. Importantly, MoLRG differs from mixture of Gaussian (MoG) distributions in diffusion models; the use of MoLRG generally focuses on learning the covariance matrices, whereas prior studies of generic MoG have primarily emphasized learning the means. 

% If we consider the denoising autoencoder (DAE) formulation for training diffusion models, then the training loss can be written as
% \begin{align*}
% \min_{\bm \theta}\ \ell(\bm \theta) := \sum_{i=1}^N \int_0^1 \lambda_t \mathbb E_{\bm \epsilon \sim \mathcal{N}(\bm 0, \bm I_n)} \left[\left\| \bm x_{\bm \theta}( s_t \bm x^{(i)} + \gamma_t \bm \epsilon, t) -  \bm x^{(i)} \right\|^2\right] \mathrm{d}t,
% \end{align*}
% Here, $s_t$ and $\gamma_t$ are the predefined hyperparameters that control the added noise schedule, and $\lambda_t$ is a weight function for the training loss. If we parameterize the DAE according to the form of the optimal posterior estimator of the above MoLRG distribution—where the subspace 
%  is learnable—we can then show that the training loss is equivalent to the canonical subspace clustering problem [8]:

\section{Summary and Notes}
\paragraph{Key Messages.}
In this Chapter, we have presented that, to pursue a general low-dimensional distribution, a universal computational approach is to start with a generic distribution and compress it towards the distribution, through (iterative) noise or entropy reduction. As we have discussed in the preceding chapter, this is even the case for arguably the simplest distributions of a single subspace, say computing PCA via the power iteration algorithm. Conceptually, the fact that deep learning is mainly conducting compression, to a large extent, explains why modern large deep models {\em do not overfit the noisy data}. As we will see in future chapters, the fundamental roles of modern deep networks are not to fit arbitrary (noisy) data distributions or functions, instead they are to realize or approximate those compressing or denoising operators! We will derive precise forms of those operators and the associated deep architectures in future chapters. Hence it becomes a simple statistical fact that, through training, these operators only become increasingly more accurate as more (noisy) training data are used. 

\paragraph{More Extensions and References.}
The use of denoising and diffusion for sampling has a rich history. The first work which is clearly about a diffusion model is probably \cite{Sohl-Dickstein2015}, but before this there are many works about denoising as a computational and statistical problem. The most relevant of these is probably \cite{hyvarinen05a}, which explicitly uses the score function to denoise (as well as perform independent component analysis). The most popular follow-ups are basically co-occurring: \cite{ho2020denoising,song2019}. Since then, thousands of papers have built on diffusion models; we will revisit this topic in \Cref{ch:autoencoding}.

Many of these works use a different stochastic process than the simple linear combination \eqref{eq:gen_additive_gaussian_noise_model}. In fact, all works listed above emphasize the need to add \textit{independent} Gaussian noise at the beginning of each step of the forward process. Theoretically-minded work actually uses Brownian motion or stochastic differential equations to formulate the forward process \cite{song2020score}. However, since linear combinations of Gaussians still result in Gaussians, the \textit{marginal distributions} of such processes still take the form of \eqref{eq:gen_additive_gaussian_noise_model}. Most of our discussion requires only that the marginal distributions are what they are, and hence our overly simplistic model is actually quite enough for almost everything. In fact, the only time where marginal distributions are not enough is when we derive an expression for \(\Ex[\vx_{s} \mid \vx_{t}]\) in terms of \(\Ex[\vx \mid \vx_{t}]\). Different (noising) processes give different such expressions, which can be used for sampling (and of course there are other ways to derive efficient samplers, such as the ever-popular DDPM sampler). The process in \eqref{eq:gen_additive_gaussian_noise_model} is a bona fide stochastic process, however, whose ``natural'' denoising iteration takes the form of the popular DDIM algorithm \cite{song2020denoising}. (Even this equivalence is not trivial; we cite \cite{de2025distributional} as a justification.) 


On top of the theoretical work \citep{li2024d} covered in \Cref{sub:sampling_denoising}, and the lineage of work that it builds on, which studies the \textit{sampling} efficiency of diffusion models when the data has low-dimensional structure, there is a large body of work which studies the \textit{training} efficiency of diffusion models when the data has low-dimensional structure. Specifically, \citet{chen2023score} and \citep{wang2024diffusion} characterized the approximation and estimation error of denoisers when the data belongs to a mixture of low-rank Gaussians, showing that the number of training samples required to accurately learn the distribution scales with the intrinsic dimension of the data rather than the ambient distribution. There is considerable \textit{methodological} work which attempts to utilize the low-dimensional structure of the data in order to do various things with diffusion models. We highlight three here: image editing \citep{chen2024exploring}, watermarking \citep{li2024shallow}, and unlearning \citep{chen2025dual}, though as always this is an inexhaustive list.


\section{Exercises and Extensions}

\begin{exercise}
    Please show that \eqref{eq:optimal_denoiser} is the optimal solution of Problem \eqref{eq:denoising_loss}. 
\end{exercise}

\begin{exercise}\label{exercise:conditional_gaussian}
  Consider random vectors $\vx \in \bR^D$ and $\vy \in \bR^d$, such that the
  pair $(\vx, \vy) \in \bR^{D + d}$ is jointly Gaussian. This means that
  \begin{equation*}
    \begin{bmatrix}
      \vx \\
      \vy
    \end{bmatrix}
    \sim
    \cN \left(
      \begin{bmatrix}
        \vmu_{\vx} \\
        \vmu_{\vy}
      \end{bmatrix}
      ,
      \begin{bmatrix}
        \vSigma_{\vx} & \vSigma_{\vx\vy} \\
        \vSigma_{\vx\vy}^\top & \vSigma_{\vy}
      \end{bmatrix}
    \right),
  \end{equation*}
  where the mean and covariance parameters are given by
  \begin{equation*}
    \vmu_{\vx} = \bE[\vx],\quad \vmu_{\vy} = \bE[\vy],\quad
    \begin{bmatrix}
      \vSigma_{\vx} & \vSigma_{\vx\vy} \\
      \vSigma_{\vx\vy}^\top & \vSigma_{\vy}
    \end{bmatrix}
    =
    \bE\left[
      \begin{bmatrix}
        \vx - \bE[\vx] \\
        \vy - \bE[\vy]
      \end{bmatrix}
      \begin{bmatrix}
        \vx - \bE[\vx] \\
        \vy - \bE[\vy]
      \end{bmatrix}^\top
      \right]
  \end{equation*}
  Assume that $\vSigma_{\vy}$ is positive definite (hence invertible); then
  positive semidefiniteness of the covariance matrix is equivalent to the Schur
  complement condition $\vSigma_{\vx} - \vSigma_{\vx\vy} \vSigma_{\vy}^{-1}
  \vSigma_{\vx\vy}^\top \succeq \Zero$.

  In this exercise, we will prove that the conditional distribution $p_{\vx \mid
  \vy}$ is Gaussian: namely,
  \begin{equation}\label{eq:gaussian-conditional-eqn}
    p_{\vx \mid \vy} \sim \cN\left(
      \vmu_{\vx} + \vSigma_{\vx\vy} \vSigma_{\vy}^{-1} (\vy - \vmu_{\vy}),
      \vSigma_{\vx} - \vSigma_{\vx\vy} \vSigma_{\vy}^{-1}
      \vSigma_{\vx\vy}^{\top}
    \right).
  \end{equation}
  A direct path to prove this result manipulates the defining ratio of
  densities $p_{\vx, \vy} / p_{\vy}$. We sketch an algebraically-concise
  argument of this form below.

  \begin{enumerate}
    \item Verify the following matrix identity for the covariance:
      \begin{equation}\label{eq:gaussian-conditional-covariance-block}
        \begin{bmatrix}
          \vSigma_{\vx} & \vSigma_{\vx\vy} \\
          \vSigma_{\vx\vy}^\top & \vSigma_{\vy}
        \end{bmatrix}
        =
        \begin{bmatrix}
          \vI_D & \vSigma_{\vx\vy}\vSigma_{\vy}^{-1} \\
          \Zero & \vI_d
        \end{bmatrix}
        \begin{bmatrix}
          \vSigma_{\vx} - \vSigma_{\vx\vy} \vSigma_{\vy}^{-1}
          \vSigma_{\vx\vy}^{\top} & \Zero \\
          \Zero & \vSigma_{\vy}
        \end{bmatrix}
        \begin{bmatrix}
          \vI_D & \Zero\\
          \vSigma_{\vy}^{-1}\vSigma_{\vx\vy}^\top & \vI_d
        \end{bmatrix}.
      \end{equation}
      One arrives at this identity by performing two rounds of (block) Gaussian
      elimination on the covariance matrix.
    \item Based on the previous identity, show that
      \begin{equation}\label{eq:gaussian-conditional-covariance-inverse-block}
        \begin{bmatrix}
          \vSigma_{\vx} & \vSigma_{\vx\vy} \\
          \vSigma_{\vx\vy}^\top & \vSigma_{\vy}
        \end{bmatrix}^{-1}
        =
		\vP
        \begin{bmatrix}
          \left(\vSigma_{\vx} - \vSigma_{\vx\vy} \vSigma_{\vy}^{-1}
          \vSigma_{\vx\vy}^{\top}\right)^{-1} & \Zero \\
          \Zero & \vSigma_{\vy}^{-1}
        \end{bmatrix}
		\vP^\top,
      \end{equation}
	  where
	  \begin{equation}
		  \vP = 
        \begin{bmatrix}
          \vI_D & \Zero\\
          -\vSigma_{\vy}^{-1}\vSigma_{\vx\vy}^\top & \vI_d
        \end{bmatrix},
	  \end{equation}
      and whenever the relevant inverses are defined.\footnote{In cases where the
      Schur complement term is not invertible, the same result holds with its
      inverse replaced by the Moore-Penrose pseudoinverse. In particular, the
      conditional distribution \eqref{eq:gaussian-conditional-eqn} becomes
      a degenerate Gaussian distribution.}
      Conclude that
      \begin{align}
        &\begin{bmatrix}
          \vx-\vmu_{\vx} \\
          \vy-\vmu_{\vy}
        \end{bmatrix}^{\top}
        \begin{bmatrix}
          \vSigma_{\vx} & \vSigma_{\vx\vy} \\
          \vSigma_{\vx\vy}^\top & \vSigma_{\vy}
        \end{bmatrix}^{-1}
        \begin{bmatrix}
          \vx-\vmu_{\vx} \\
          \vy-\vmu_{\vy}
        \end{bmatrix}
        \\
        &\qquad=
        \begin{bmatrix}
          \vx - \vm \\
          \vy - \vmu_{\vy}
        \end{bmatrix}^\top
        \begin{bmatrix}
          \left(\vSigma_{\vx} - \vSigma_{\vx\vy} \vSigma_{\vy}^{-1}
          \vSigma_{\vx\vy}^{\top}\right)^{-1} & \Zero \\
          \Zero & \vSigma_{\vy}^{-1}
        \end{bmatrix}
        \begin{bmatrix}
          \vx - \vm \\
          \vy - \vmu_{\vy}
        \end{bmatrix},
      \end{align}
	  where $\vm = (\vmu_{\vx} + \vSigma_{\vx\vy}\vSigma_{\vy}^{-1}(\vy
	  - \vmu_{\vy})$.
      %(\textit{Hint: To economize algebraic manipulations, note that the first
      %and last matrices on the RHS of
      %\Cref{eq:gaussian-conditional-covariance-block} are transposes of one
      %another.})
    \item By dividing $p_{\vx, \vy} / p_{\vy}$, prove
      \Cref{eq:gaussian-conditional-eqn}. (\textit{Hint: Using the previous
      identities, only minimal algebra should be necessary. For the normalizing
      constant, use \Cref{eq:gaussian-conditional-covariance-inverse-block} to
      factor the determinant similarly.})

  \end{enumerate}


\end{exercise}

\begin{exercise}\label{exercise:sherman_morrison_woodbury_identity}
    Show the Sherman-Morrison-Woodbury identity, i.e., for matrices \(\vA\), \(\vC\), \(\vU\), \(\vV\) such that \(\vA\), \(\vC\), and \(\vA + \vU\vC\vV\) are invertible,
    \begin{equation}
        (\vA + \vU\vC\vV)^{-1} = \vA^{-1} - \vA^{-1}\vU(\vC^{-1} + \vV\vA^{-1}\vU)^{-1}\vV\vA^{-1}
    \end{equation}
\end{exercise}

\begin{exercise}\label{exercise:generalizing_results_to_different_noise_models}
    Rederive the following, assuming \(\vx_{t}\) follows the generalized noise model \eqref{eq:gen_additive_gaussian_noise_model}.
    \begin{itemize}
        \item Tweedie's formula: \eqref{eq:gen_tweedie}.
        \item The DDIM iteration: \eqref{eq:gen_denoising_iteration}.
        \item The Bayes optimal denoiser for a Gaussian mixture model: \eqref{eq:gen_gmm_bayes_optimal_denoiser}.
    \end{itemize}
\end{exercise}

\begin{exercise}\label{exercise:implement_denoising_processes}
\begin{enumerate}
    \item Implement the formulae derived in \Cref{exercise:generalizing_results_to_different_noise_models}, building a sampler for Gaussian mixtures.
    \item Reproduce \Cref{fig:ve_forward_denoising} and \Cref{fig:vp_gmm_denoising}.
    \item We now introduce a separate process called \textit{Flow Matching (FM)}, as follows:
    \begin{equation}
        \alpha_{t} = 1 - t, \qquad \sigma_{t} = t.
    \end{equation}
    Implement this process using the same framework, and test it for sampling in high dimensions. Which process seems to give better or more stable results?
\end{enumerate}
\end{exercise}




 

\end{document}
